% Options for packages loaded elsewhere
\PassOptionsToPackage{unicode}{hyperref}
\PassOptionsToPackage{hyphens}{url}
\PassOptionsToPackage{dvipsnames,svgnames,x11names}{xcolor}
%
\documentclass[
  letterpaper,
  DIV=11,
  numbers=noendperiod]{scrartcl}

\usepackage{amsmath,amssymb}
\usepackage{iftex}
\ifPDFTeX
  \usepackage[T1]{fontenc}
  \usepackage[utf8]{inputenc}
  \usepackage{textcomp} % provide euro and other symbols
\else % if luatex or xetex
  \usepackage{unicode-math}
  \defaultfontfeatures{Scale=MatchLowercase}
  \defaultfontfeatures[\rmfamily]{Ligatures=TeX,Scale=1}
\fi
\usepackage{lmodern}
\ifPDFTeX\else  
    % xetex/luatex font selection
  \setmainfont[]{Times New Roman}
\fi
% Use upquote if available, for straight quotes in verbatim environments
\IfFileExists{upquote.sty}{\usepackage{upquote}}{}
\IfFileExists{microtype.sty}{% use microtype if available
  \usepackage[]{microtype}
  \UseMicrotypeSet[protrusion]{basicmath} % disable protrusion for tt fonts
}{}
\usepackage{xcolor}
\setlength{\emergencystretch}{3em} % prevent overfull lines
\setcounter{secnumdepth}{-\maxdimen} % remove section numbering
% Make \paragraph and \subparagraph free-standing
\ifx\paragraph\undefined\else
  \let\oldparagraph\paragraph
  \renewcommand{\paragraph}[1]{\oldparagraph{#1}\mbox{}}
\fi
\ifx\subparagraph\undefined\else
  \let\oldsubparagraph\subparagraph
  \renewcommand{\subparagraph}[1]{\oldsubparagraph{#1}\mbox{}}
\fi


\providecommand{\tightlist}{%
  \setlength{\itemsep}{0pt}\setlength{\parskip}{0pt}}
% longtable/footnotehyper/etoolbox longtable-patch removed: this
% main-text-only file contains no longtable environments (those are all
% in the Supporting Information file).
\usepackage{booktabs,array}
\usepackage{calc} % for calculating minipage widths
\usepackage{graphicx}
\makeatletter
\def\maxwidth{\ifdim\Gin@nat@width>\linewidth\linewidth\else\Gin@nat@width\fi}
\def\maxheight{\ifdim\Gin@nat@height>\textheight\textheight\else\Gin@nat@height\fi}
\makeatother
% Scale images if necessary, so that they will not overflow the page
% margins by default, and it is still possible to overwrite the defaults
% using explicit options in \includegraphics[width, height, ...]{}
\setkeys{Gin}{width=\maxwidth,height=\maxheight,keepaspectratio}
% Set default figure placement to htbp
\makeatletter
\def\fps@figure{htbp}
\makeatother
% definitions for citeproc citations
\NewDocumentCommand\citeproctext{}{}
\NewDocumentCommand\citeproc{mm}{%
  \begingroup\def\citeproctext{#2}\cite{#1}\endgroup}
\makeatletter
 % allow citations to break across lines
 \let\@cite@ofmt\@firstofone
 % avoid brackets around text for \cite:
 \def\@biblabel#1{}
 \def\@cite#1#2{{#1\if@tempswa , #2\fi}}
\makeatother
\newlength{\cslhangindent}
\setlength{\cslhangindent}{1.5em}
\newlength{\csllabelwidth}
\setlength{\csllabelwidth}{3em}
\newenvironment{CSLReferences}[2] % #1 hanging-indent, #2 entry-spacing
 {\begin{list}{}{%
  \setlength{\itemindent}{0pt}
  \setlength{\leftmargin}{0pt}
  \setlength{\parsep}{0pt}
  % turn on hanging indent if param 1 is 1
  \ifodd #1
   \setlength{\leftmargin}{\cslhangindent}
   \setlength{\itemindent}{-1\cslhangindent}
  \fi
  % set entry spacing
  \setlength{\itemsep}{#2\baselineskip}}}
 {\end{list}}
\usepackage{calc}
\newcommand{\CSLBlock}[1]{\hfill\break\parbox[t]{\linewidth}{\strut\ignorespaces#1\strut}}
\newcommand{\CSLLeftMargin}[1]{\parbox[t]{\csllabelwidth}{\strut#1\strut}}
\newcommand{\CSLRightInline}[1]{\parbox[t]{\linewidth - \csllabelwidth}{\strut#1\strut}}
\newcommand{\CSLIndent}[1]{\hspace{\cslhangindent}#1}

\usepackage{booktabs}
\usepackage{array}
\usepackage{float}
\usepackage{colortbl}
\usepackage{xcolor}
\KOMAoption{captions}{tableheading}
\usepackage{hyperref}
\addtokomafont{disposition}{\rmfamily}
\usepackage{siunitx}
\newcommand{\ca}{$C_\mathrm{a}$}
\newcommand{\caequals}[1]{$C_\mathrm{a} = \SI{#1}{\micro\mol\raiseto{-1}\mol}$}
\newcommand{\fgmax}{$f_\mathrm{gmax}$}
\newcommand{\gcl}{$l_\mathrm{gc}$}
\newcommand{\gi}{$g_\mathrm{i}$}
\newcommand{\gf}{$g_\mathrm{f}$}
\newcommand{\gmax}{$g_\mathrm{max}$}
\newcommand{\gop}{$g_\mathrm{op}$}
\newcommand{\gsw}{$g_\mathrm{sw}$}
\newcommand{\loggi}{$\log(g_\mathrm{i})$}
\newcommand{\loggmax}{$\log(g_\mathrm{max})$}
\newcommand{\ppfd}{$\mathrm{PPFD}$}
\newcommand{\ppfdequals}[1]{$\mathrm{PPFD} = \SI{#1}{\micro\mole\per\meter\squared\per\second}$}
\newcommand{\ppfdqty}[1]{$\SI{#1}{\micro\mole\per\meter\squared\per\second}$}
\newcommand{\rh}{$\mathrm{RH}$}
\newcommand{\rhequals}[1]{$\mathrm{RH} = #1\%$}
\newcommand{\tleafequals}[1]{$T_\mathrm{leaf} = \SI{#1}{\degreeCelsius}$}
\usepackage{setspace}
% The physics package is unavailable on the publisher's LaTeX system;
% define the one macro (\dv, total derivative) actually used below
% instead of loading the package.
\newcommand{\dv}[2]{\frac{\mathrm{d}#1}{\mathrm{d}#2}}
\DeclareMathOperator{\logit}{logit}
\makeatletter
\@ifpackageloaded{caption}{}{\usepackage{caption}}
\AtBeginDocument{%
\ifdefined\contentsname
  \renewcommand*\contentsname{Table of contents}
\else
  \newcommand\contentsname{Table of contents}
\fi
\ifdefined\listfigurename
  \renewcommand*\listfigurename{List of Figures}
\else
  \newcommand\listfigurename{List of Figures}
\fi
\ifdefined\listtablename
  \renewcommand*\listtablename{List of Tables}
\else
  \newcommand\listtablename{List of Tables}
\fi
\ifdefined\figurename
  \renewcommand*\figurename{Figure}
\else
  \newcommand\figurename{Figure}
\fi
\ifdefined\tablename
  \renewcommand*\tablename{Table}
\else
  \newcommand\tablename{Table}
\fi
}
\@ifpackageloaded{float}{}{\usepackage{float}}
\floatstyle{ruled}
\@ifundefined{c@chapter}{\newfloat{codelisting}{h}{lop}}{\newfloat{codelisting}{h}{lop}[chapter]}
\floatname{codelisting}{Listing}
\newcommand*\listoflistings{\listof{codelisting}{List of Listings}}
\makeatother
\makeatletter
\makeatother
\makeatletter
\@ifpackageloaded{caption}{}{\usepackage{caption}}
\@ifpackageloaded{subcaption}{}{\usepackage{subcaption}}
\makeatother
\usepackage{bookmark}

\urlstyle{same} % disable monospaced font for URLs
\hypersetup{
  pdftitle={Guard cell size and initial conductance influence stomatal closure kinetics},
  pdfauthor={Christopher D. Muir\^{}\{1,2*\}; Wei Shen Lim\^{}\{1\}},
  colorlinks=true,
  linkcolor={blue},
  filecolor={Maroon},
  citecolor={Blue},
  urlcolor={Blue},
  pdfcreator={LaTeX via pandoc}}

\title{Guard cell size and initial conductance influence stomatal
closure kinetics}
\author{Christopher D. Muir\(^{1,2*}\) \and Wei Shen Lim\(^{1}\)}
\date{}

\begin{document}
\maketitle

\begin{center}
$^1$School of Life Sciences, University of Hawaiʻi; Mānoa, Hawaiʻi, USA\\
$^2$Department of Biology, University of Wisconsin, Madison; Wisconsin, USA\\
\medskip
*Corresponding author: cdmuir@wisc.edu\\
\medskip
ORCID: Christopher D. Muir — \href{https://orcid.org/0000-0003-2555-3878}{0000-0003-2555-3878}\\
ORCID: Wei Shen Lim — \href{https://orcid.org/0009-0001-1406-7957}{0009-0001-1406-7957}\\
\medskip
Keywords: guard cell size, leaf gas exchange, phylogenetic heritability, \textit{Solanum}, stomatal kinetics, vapor pressure deficit
\end{center}

\begin{flushleft}
Running head: Guard cell size, initial conductance, and stomatal kinetics
\end{flushleft}

% \linenumbers
% \onehalfspacing

\newpage{}

\subsection{Summary}\label{summary}

\begin{itemize}
\item
  In fluctuating environments, the kinetics of stomatal opening and
  closing influence the balance between carbon gain and water loss.
  Smaller guard cells may respond faster to fluctuating environmental
  conditions because of their greater surface area for osmolyte flux
  relative to cell volume. A related hypothesis is that operational
  stomatal conductance (\gop) is often well below its theoretical
  maximum (\gmax) because at this stomatal aperture, guard cell volume
  is poised to change rapidly with small changes in turgor pressure.
\item
  We analyzed 2,125 estimates of stomatal closure kinetics in response
  to an abrupt increase in vapor pressure deficit (VPD) among 29 diverse
  wild tomato populations in the genus \emph{Solanum}.
\item
  Leaves with small guard cells and a lower initial stomatal conductance
  (\gi) closed faster, but each explained variation in kinetic
  parameters at different levels of biological organization. Guard cell
  size had high phylogenetic heritability and varied relatively little
  within populations, whereas \gi{} varied mostly among individuals and
  between light intensity treatments.
\item
  Smaller stomata can be speedier, but the current physiological state
  of the leaf also affects the responsiveness of stomatal conductance to
  fluctuations. Selection on stomatal speed may influence not only
  anatomical traits like guard cell size, but also physiological
  controls on \gop.
\end{itemize}

\newpage{}

\subsection{Introduction}\label{introduction}

In nature, change is the only constant. One way that vascular plants
(tracheophytes) cope with change is by adjusting stomatal pore aperture
to optimize the trade-off between carbon gain and water loss (Cowan \&
Farquhar, 1977; Raven, 2002; Hetherington \& Woodward, 2003). In the
absence of physical limits and energetic tradeoffs, natural selection
would favor plants that could instantaneously adjust stomatal
conductance (\gsw) to perfectly track dynamic environmental conditions
(Lawson \emph{et al.}, 2010). In reality, stomatal responses take time,
creating a lag between actual and optimal stomatal aperture. For
example, many species use hydroactive feedback mechanisms to precisely
control stomatal responses to changing vapor pressure deficit (VPD), the
driving force for evaporation (Buckley, 2019), but hydroactive control
takes time. Leaves must sense a change in VPD through as-yet-unknown
mechanisms and then transduce a signal that causes efflux of osmotic
solutes to alter guard cell turgor more than the surrounding epidermal
cells. Abscisic acid (ABA) dependent responses require several minutes
to accumulate or synthesize hormone pools \emph{de novo} (Xie \emph{et
al.}, 2006; Bauer \emph{et al.}, 2013; McAdam \emph{et al.}, 2016; Desai
\& Stroock, 2025; Pichaco \emph{et al.}, 2026); ABA-independent stomatal
responses require changes in gene expression and protein synthesis that
take time (Jalakas \emph{et al.}, 2021). Once anion channels are
activated, the rate of change in stomatal aperture depends on the rate
of change in guard cell turgor (\(d P_\mathrm{g} / dt\)) and the
sensitivity of aperture to change in turgor
(\(d a_\mathrm{S} / d P_\mathrm{g}\)). As we discuss below, these two
factors (\(d P_\mathrm{g} / dt\) and
\(d a_\mathrm{S} / d P_\mathrm{g}\)) may be influenced by guard cell
size and the ratio of operational to maximum stomatal conductance
(\fgmax), respectively.

Over a century ago Darwin (1898) observed that after an initial
increase, leaf transpiration decreases in response to decreasing
humidity. He correctly attributed changes in transpiration to changes in
stomatal aperture and subsequent studies have confirmed that stomatal
aperture responds to humidity (Lange \emph{et al.}, 1971) and, more
specifically, transpiration rate (Mott \& Parkhurst, 1991; Monteith,
1995), which is the product of VPD and \gsw. The VPD is related to
relative humidity (RH) but also temperature. In angiosperms and some
pteridophytes, a sudden increase in VPD, usually induced by lower RH at
a constant temperature, causes a transient increase in \gsw, as Darwin
observed (Cowan, 1972; Buckley \emph{et al.}, 2011; Westbrook \& McAdam,
2021). After a period of several minutes, stomata close until the leaf
reaches a new steady-state \gsw. These phases of stomatal response to
VPD are often referred to as the ``wrong-way response'' (WWR) and
``right-way response'' (RWR), respectively. The normative `wrong' and
`right' appellations refer to the presumed adaptive significance. At
steady state, \gsw{} is negatively correlated with VPD, all else being
equal, which is interpreted as an adaptive response to the marginal
carbon cost of water (Givnish \& Vermeij, 1976; Cowan \& Farquhar, 1977;
Buckley \emph{et al.}, 2017b). Models built around this assumption
generally predict steady state \gsw{} well, although the functional form
and presumed fitness costs are areas of active research (Damour \emph{et
al.}, 2010; Medlyn \emph{et al.}, 2011; Prentice \emph{et al.}, 2014;
Sperry \emph{et al.}, 2017; Potkay \emph{et al.}, 2025).

Stomatal responses are as much about the journey (kinetics) as they are
about the destination (steady state). If stomatal response kinetics are
too slow relative to the temporal grain size of environmental
variability, then it would be futile for \gsw{} to track environmental
variation since the leaf will constantly lag the environment,
perpetually fighting the last war. This may explain why \gsw{} can be
insensitive to short sun/shade flecks (Knapp \& Smith, 1990) and shaded
leaves are likely primed to cope with the hydraulic demands of sudden
increases in light, temperature, and VPD (Schymanski \emph{et al.},
2013). If stomatal responses can track environmental variation, then
minimizing the lag time can greatly improve resource-use efficiency
(Lawson \emph{et al.}, 2010). For example, optogenetic control of
stomatal aperture in response to fluctuating light increases water-use
efficiency in \emph{Arabidopsis} compared to wild-type responses
(Papanatsiou \emph{et al.}, 2019). Therefore the rate of stomatal
opening and closing is predicted to have a major effect on how natural
selection optimizes stomatal control in variable environments (Buckley
\emph{et al.}, 2023).

Much of the research on stomatal kinetics focuses on light responses
because light intensity is usually the most important, rapidly
fluctuating environmental variable affecting photosynthesis. In nature,
fluctuating light intensity simultaneously alters leaf temperature and
VPD (e.g. Ishida \emph{et al.}, 1992). We focus on stomatal closure
kinetics in response to an abrupt increase in VPD, rather than light,
for four reasons. First of all, the primary aims and \emph{a priori}
hypotheses of this study were not to study stomatal kinetic parameters
and are described in Muir \emph{et al.} (2025). The idea to
opportunistically use \gsw-response curves to study stomatal kinetics
came after the data had been collected, but not yet analyzed for this
purpose. Second, VPD responses exhibit a similar functional form as
light-responses, suggesting similar underlying mechanisms (Buckley,
2019). Factors that affect stomatal response kinetics to VPD likely
generalize to other environmental variables. Third, humidity responses
tend to dominate over and are faster than other environmental responses
(Aasamaa \& Sõber, 2011a,b), although there is variation among species
(Merilo \emph{et al.}, 2014). Thus, factors that constrain maximum rates
of stomatal closure are most likely to be observed in response to VPD.
Finally, the ecological salience of VPD responses is increasing as
climate change warms the atmosphere (Grossiord \emph{et al.}, 2020).

Fully describing stomatal responses to the environment is complex
(Buckley \emph{et al.}, 2003; Jezek \emph{et al.}, 2021; Desai \&
Stroock, 2025), but the dynamics of \gsw{} in response to an abrupt
environmental change is well described by relatively simple
phenomenological functions (Vico \emph{et al.}, 2011; McAusland \emph{et
al.}, 2016). A four-parameter equation describes stomatal response
kinetics to a step change in light intensity across dozens of plant
species (Woning \emph{et al.}, 2026):

\begin{equation}\phantomsection\label{eq-weibull}{
  g_\mathrm{sw,t} = g_\mathrm{f} + (g_\mathrm{i} - g_\mathrm{f}) e^{-\left(\frac{t}{\tau}\right)^\lambda}.
}\end{equation}

In this equation, \(g_\mathrm{sw,t}\) is stomatal conductance at time
\(t\), \gi{} is the initial \gsw{} before the step change, \gf{} is the
final \gsw{} after the step change, \(\tau\) is a time constant that
describes how quickly stomata respond to the step change, and
\(\lambda\) is a shape parameter that describes how the response rate
changes over time. When \(\lambda = 1\), the response rate is constant
over time and \(\tau\) represents the time required for \gsw{} to reach
63.2\% (\(1 - e^{-1}\)) of the way between initial and final values.
When \(\lambda > 1\), the response rate increases over time, and when
\(\lambda < 1\), the response rate decreases over time. Woning \emph{et
al.} (2026) refer to \(\lambda\) as a lag-time parameter because
empirically \(\lambda > 1\). We refer to stomatal kinetic parameters
\(\tau\) (time constant) and \(\lambda\) (lag time) henceforth. The
first goal of this study is to evaluate whether kinetics of the RWR
after a step change in VPD are adequately described by
Equation~\ref{eq-weibull}. If so, it suggests that some fundamentally
similar mechanisms are shared between light and hydroactive stomatal
responses (Grantz \& Zeiger, 1986).

The main goal of this study is to test hypotheses about which stomatal
anatomical traits affect kinetic parameters \(\tau\) and \(\lambda\). We
consider two interconnected hypotheses that have thus far been treated
in isolation. The first hypothesis is that smaller guard cells open and
close faster because of their intrinsically greater surface area to
volume ratio (Franks \& Farquhar, 2007; Drake \emph{et al.}, 2013;
Lawson \& Vialet-Chabrand, 2019; Woning \emph{et al.}, 2026). For
approximately cylindrical cell geometries, like that found in guard
cells, surface area increases linearly with radius, whereas volume
increases in proportion to the radius squared. Consequently, larger
guard cells will require more time, all else being equal, to achieve a
given change in turgor pressure and, hence, stomatal pore aperture. A
plant with large guard cells can partially compensate for this geometric
constraint by increasing the activity and/or density of ion and solute
transporters in the guard cell plasma membrane (Homann \& Thiel, 2002;
Lawson \& Blatt, 2014; Raven, 2014). However, compensation may be
limited by the available membrane area and/or energetic costs of
maintaining high transporter densities (Vico \emph{et al.}, 2011).
Empirical support for the hypothesis that smaller guard cells are faster
is mixed and has primarily come from step changes in light (Drake
\emph{et al.}, 2013; Elliott-Kingston \emph{et al.}, 2016; McAusland
\emph{et al.}, 2016; Kardiman \& Ræbild, 2018; Haworth \emph{et al.},
2023; Yoshiyama \emph{et al.}, 2024; Brench \emph{et al.}, 2026; Tang
\emph{et al.}, 2026). Meta-analysis of dozens of species suggests that
stomatal size alone is weakly related to kinetics and modulated by guard
cell shape (Woning \emph{et al.}, 2026). The relationship between guard
cell size and stomatal response kinetics to leaf water potential
(Aasamaa \emph{et al.}, 2001) and VPD (Kübarsepp \emph{et al.}, 2020;
Qie \emph{et al.}, 2024) is less well studied, but similarly equivocal.

One factor that might complicate the relationship between size and speed
is aperture. Stomatal aperture can vary from near 0 when guard cell
turgor pressure is low and asymptotically approach \gmax{} as guard cell
turgor pressure approaches \(\infty\) (Figure~\ref{fig-conceptual}d).
The relationship between guard cell turgor pressure and stomatal
aperture is nonlinear (Franks \& Farquhar, 2001) as cell walls become
less elastic at greater turgor (Franks \emph{et al.}, 2001). The
nonlinear relationship between guard cell turgor and pore aperture
implies that if the rate of turgor change is constant
(\(\dv{P}{t} = C\)), the rate of change in aperture, and hence stomatal
conductance, will vary depending on the initial aperture (\(a_0\))
relative to the maximum aperture (\(a_\mathrm{max}\)). Hence,
\(\dv{a}{t} = f(a_0 / a_\mathrm{max})\). When initial aperture is high,
\gsw{} will respond more slowly than when initial aperture is low. Since
we generally do not observe individual stomatal aperture, relating this
idea to leaf-level \gsw{} requires scaling by stomatal density,
biophysical constants, and morphological constants (e.g. Sack \&
Buckley, 2016). We will use the term ``fraction of anatomical maximum
stomatal conductance'', symbolized as \fgmax{} (this is denoted as
\(g_\mathrm{ratio}\) by Xie \emph{et al.} (2022)). This value should be
proportional to the average aperture divided by its maximum aperture for
any arbitrary stomatal density. Leaves operating at \fgmax{} closer to
unity will, all else being equal, be slower to respond than leaves
operating with a \fgmax{} close to zero, independent of stomatal
density. Franks \emph{et al.} (2012) hypothesized that selection on
stomatal density would maintain typical operational stomatal conductance
(\gop) relative to \gmax{} at a sufficiently low value that \gsw{} would
be sensitive to relatively small changes in guard cell turgor pressure.
Consistent with this hypothesis, the \gop:\gmax{} ratio is often near
0.25 (McElwain \emph{et al.}, 2016; Murray \emph{et al.}, 2020), well
below \(f_\mathrm{gmax} = 1\), and in a range where \gsw{} would be
responsive to small changes in guard cell turgor pressure. There is
likely variation within and among species and growth forms in the
\gop:\gmax{} ratio (Xie \emph{et al.}, 2022; Ochoa \emph{et al.}, 2024).

\begin{figure}

\centering{

\Large{[Figure 1 goes here]}

}

\caption{\label{fig-conceptual}\textbf{Conceptual overview of the two hypotheses linking stomatal anatomy to closure kinetics.}
\textbf{a.} Stomatal conductance (\gsw) in response to a step increase
in vapor pressure deficit (VPD) typically exhibits a transient wrong-way
response (WWR), in which \gsw{} briefly rises, followed by the right-way
response (RWR), in which \gsw{} declines to a lower steady state. The
RWR is well described by a Weibull function (dashed line;
Equation~\ref{eq-weibull}) with initial conductance \(g_\mathrm{i}\),
final conductance \(g_\mathrm{f}\), and kinetic parameters \(\tau\) and
\(\lambda\). \textbf{b.} The time constant \(\tau\) sets the overall
speed of closure, and the lag time \(\lambda\) determines whether the
response rate is roughly constant (\(\lambda \approx 1\); solid) or
accelerates over time (\(\lambda > 1\); dashed). \textbf{c.} Smaller
guard cells have a greater surface area to volume ratio (SA:V; schematic
cross-sections shown at right) and are predicted to exchange osmolytes
more quickly across the plasma membrane, yielding a faster change in
guard cell turgor pressure and faster closure. \textbf{d.} The
relationship between guard cell turgor pressure (\(P_g\)) and \gsw{} is
nonlinear and saturating because guard cell walls become less elastic at
higher turgor (Franks \& Farquhar, 2001; Franks \emph{et al.}, 2001).
Consequently, the same change in turgor pressure produces a larger
change in \gsw{} when stomata are at low \fgmax{} (orange; steep tangent
slope) than at high \fgmax{} (blue; shallow tangent slope), predicting
slower closure when \fgmax{} is high.}

\end{figure}%

Putting these two hypotheses together, we predict that both guard cell
size and \fgmax{} will influence stomatal kinetics, but possibly at
different scales of biological organization. Guard cell size tends to
vary less than stomatal density or aperture at many biological scales.
For example, on a mature leaf, guard cell length is nearly constant even
as cell volume and pore aperture change with turgor (Willmer \& Fricker,
1996; Franks \emph{et al.}, 2001). Compared to stomatal density, guard
cell size can be less variable (Lammertsma \emph{et al.}, 2011; Bucher
\emph{et al.}, 2016) and evolve more slowly between species (Muir
\emph{et al.}, 2023). If there is relatively little plastic or genetic
variation among individuals within a species, then most of the variance
in stomatal kinetics within species cannot be explained by variation in
guard cell size. In contrast, \fgmax{} can vary immensely within and
among individuals because stomatal aperture responds dynamically over
the day and in response to environmental variation. Because guard cell
size and aperture typically vary at different levels of biological
organization, we predict that guard cell size likely explains more
variation in stomatal kinetics among species, whereas \fgmax{} will
explain more variation within species. Within species variation can
arise from either genetic differences among individuals or environmental
plasticity. We will also test if adaxial stomata respond faster than
abaxial stomata (Haworth \emph{et al.}, 2018; Wall \emph{et al.}, 2022;
Woning \emph{et al.}, 2026) and, if so, whether this difference can be
explained by variation in guard cell size and/or \fgmax{} between
surfaces.

One challenge in determining whether traits covary at higher levels of
organization is that statistical procedures to account for phylogenetic
nonindependence can obscure ecologically relevant covariation caused by
natural selection (Westoby \emph{et al.}, 1995; Uyeda \emph{et al.},
2018). For example, in a meta-analysis of stomatal kinetic responses to
light, the explanatory power of traits like guard cell size and stomatal
ratio declined in phylogenetic compared to nonphylogenetic regression
(Woning \emph{et al.}, 2026). This finding may be caused in large part
because grasses have small stomata, high adaxial:abaxial stomatal ratio,
and rapid stomatal kinetics. Are these associations a coincidence caused
by shared ancestry or ecologically important covariation that helps
explain how grasses adapt to their environment? Recent advances in
phylogenetic comparative methods use hierarchical, multiresponse
phylogenetic models (Muir \emph{et al.}, 2023; Westoby \emph{et al.},
2023; Halliwell \emph{et al.}, 2025) to tease apart ``conservative trait
correlations'' that act at higher levels of organization (e.g.~among
species) from trait correlations that explain variation at lower levels
of organization (e.g.~among individuals within species and plastic
responses). We use this approach to test whether guard cell size and
\fgmax{} explain variation in stomatal kinetics at different levels of
biological organization. Throughout, we use guard cell length (\gcl) as
a proxy for size (Sack \& Buckley, 2016).

The four main questions we address in this study are:

\begin{enumerate}
\def\labelenumi{\arabic{enumi}.}
\tightlist
\item
  Are VPD kinetics adequately characterized by the same functional form
  in Equation~\ref{eq-weibull} as light responses?
\item
  Are \gcl{} and \fgmax{} positively associated with \(\tau\)?
\item
  At what biological levels of organization do \gcl{} and \fgmax{} vary
  and influence \(\tau\)?
\item
  Are abaxial-only stomatal responses slower than whole leaf stomatal
  responses?
\end{enumerate}

As noted above, we are opportunistically analyzing data from an
experiment designed to answer different questions. Before analyzing the
kinetic data, we planned to test whether 1) guard cell size would
influence \(\tau\) and 2) that \(\tau\) would be greater on the
abaxial-only leaves based on a recent meta-analysis of stomatal
responses to light (Woning \emph{et al.}, 2026). We made the decision to
test for an association between \fgmax{} and \(\tau\) after exploratory
analyses revealed a potential association. Adding hypotheses after some
results have been analyzed, also known as ``Hypothesizing After Results
Known'' (HARKing), can lead to reports of spurious associations that are
not repeatable in later studies (Kerr, 1998; Smaldino \& McElreath,
2016). Therefore, our conclusions about the effect of \fgmax{} on
stomatal kinetics should be interpreted with caution until they are
replicated.

\subsection{Materials and Methods}\label{sec-methods}

\subsubsection{Populations and
phylogeny}\label{populations-and-phylogeny}

We grew 29 populations of wild tomato (Table~S1),
including representatives of all described species of \emph{Solanum}
sect. \emph{Lycopersicon} and sect. \emph{Lycopersicoides} (Peralta
\emph{et al.}, 2008), as described in Muir \emph{et al.} (2025). These
populations were selected because they encompass the breadth of climatic
variation in the wild tomato clade (Pease \emph{et al.}, 2016). Due to
constraints on growth space and time, we spread out measurements over
61.1 weeks. Replicates within population were evenly spread out over
this period to prevent confounding of temporal variation in growth
conditions with variation among populations. For phylogenetic
comparative analyses, we used the RAxML whole-transcriptome concatenated
phylogeny based on 2,745 100-kb genomic windows from Pease \emph{et al.}
(2016). Two of our populations were not in this tree. We used \emph{S.
galapagense} S.C. Darwin \& Peralta accession LA1044 in place of LA3909.
\emph{S. pennellii} Correll accession LA0750 was added as sister to
LA0716. The node separating LA0750 from LA0716 was placed halfway along
the branch to the next deepest node.

\subsubsection{Plant growth conditions and light
treatments}\label{plant-growth-conditions-and-light-treatments}

A thorough description of plant growth conditions can be found in Muir
\emph{et al.} (2025), therefore we summarize the most important
information here. Seeds provided by the Tomato Genetics Resource Center
germinated on moist paper in plastic boxes after soaking for 30-60
minutes in a 50\% (volume per volume) solution of household bleach and
water, followed by a thorough rinse. We transferred seedlings to
cell-pack flats containing Pro-Mix BX potting mix (Premier Tech,
Rivière-du-Loup, Quebec, Canada) once cotyledons fully emerged,
typically within 1-2 weeks of sowing. We grew seeds and seedlings for
both sun and shade treatments under the same environmental conditions
(12:12 h, 24.3:\(\SI{21.7}{\degreeCelsius}\), 49.6:58.4 RH day:night
cycle). LED light provided
\(\mathrm{PPFD} = \SI{267}{\micro\mole\per\meter\squared\per\second}\)
(Fluence RAZRx, Austin, Texas, USA) during the germination and seedling
stages.

Immediately prior to starting growth light intensity treatments (sun and
shade), we transplanted seedlings to 3.78 L plastic pots containing 60\%
Pro-Mix BX potting mix, 20\% coral sand (Pro-Pak, Honolulu, Hawaiʻi,
USA), and 20\% cinders (Niu Nursery, Honolulu, Hawaiʻi, USA) with slow
release NPK fertilizer following manufacturer instructions (Osmocote
Smart-Release Plant Food Flower \& Vegetable, The Scotts Company,
Marysville, Ohio, USA). Percentage composition is on a volume basis. We
watered to field capacity three times per week to prevent drought
stress. Seedlings were randomly assigned in alternating order within
population to the sun or shade treatment during transplanting. The
average daytime \ppfd{} was \ppfdqty{761} and \ppfdqty{115} for sun and
shade treatments, respectively.

\subsubsection{Humidity response curves}\label{humidity-response-curves}

We selected a fully expanded, unshaded leaf at least six leaves above
the cotyledons during early vegetative growth. This typically meant that
plants had grown in light treatments for \(\approx 4\) weeks, ensuring
they had time to sense and respond developmentally to the light
intensity of the treatment rather than the seedling conditions (Schoch
\emph{et al.}, 1980).

We measured stomatal conductance over time after a step change in
humidity at a constant leaf temperature, which we refer to this as a
humidity-response curve. We measured humidity-response curves using a
portable infrared gas analyzer (LI-6800PF, LI-COR Biosciences, Lincoln,
Nebraska, USA). During gas exchange measurements, light-acclimated
plants were placed under LEDs dimmed to match their light treatment. We
collected four humidity-response curves per leaf, an amphistomatous
(untreated) curve and a pseudohypostomatous (treated) curve at high
light-intensity (\ppfdequals{2000}; 97.8:2.24 red:blue) and low
light-intensity (\ppfdequals{150}; 87.0:13.0 red:blue). We always
measured high light-intensity curves first because photosynthetic
downregulation is faster than upregulation in these species. As a
consequence, measurement light intensity is confounded with measurement
order in our design: low light-intensity curves were always measured
after high light-intensity curves on the same leaf, following a
low-humidity closure event, so we cannot fully separate an effect of
light intensity \emph{per se} from order, recovery state, or prior VPD
exposure (see Discussion). Pseudohypostomatous leaves are the same as
the amphistomatous leaves except that gas exchange through the upper
(adaxial) surface is blocked by a neutral density plastic (propafilm).
For brevity, we hereafter refer to these as ``pseudohypo'' and ``amphi''
leaf types. Comparing amphi and pseudohypo leaf types enabled us to test
whether stomatal kinetics are different for ab- and adaxial stomata on
the same leaf. To compensate for reduced transmission, we increased
incident \ppfd{} for pseudohypo leaves by a factor 1/0.91, the inverse
of the measured transmissivity of the propafilm. We also set the
stomatal conductance ratio, for purposes of calculating boundary layer
conductance, to 0 for pseudohypo leaves following manufacturer
directions. To control for order effects, we alternated between starting
with amphi or pseudohypo measurements. We made measurements over two
days. On the first day, we measured high and low light-intensity curves
for either amphi or pseudohypo leaves; on the second day, we measured
high and low light-intensity curves on the other leaf type at
approximately the same time of day.

In all humidity-response curves, we acclimated the focal leaf ambient
CO\(_2\) (\caequals{415}), \tleafequals{25.0}, high light
(\ppfdequals{2000}), and high relative humidity (\rhequals{70};
\(\mathrm{VPD} \approx \SI{0.82}{\kPa}\)) until \gsw{} reached its
maximum. After that, we scrubbed water vapor from the incoming air using
silica desiccant as quickly as possible, resulting in a water vapor
concentration in the incoming air to be on average
\(\SI{0.36}{\milli\mol\per\mol}\). This treatment induced rapid stomatal
closure, but it did not standardize chamber air humidity because leaf
transpiration introduced water vapor. Consequently, the RH was
systematically higher at high light intensity, in sun plants, and amphi
leaf types (Table~S2) because the \gsw{} was typically
greater in these conditions. In all treatment conditions, the final RH
(\(5.37\%-13.8\%\)) and VPD (\(\SI{2.72}{\kPa} - \SI{3}{\kPa}\))
elicited rapid stomatal response and is comparable to treatments used in
similar experiments (Merilo \emph{et al.}, 2014; Kübarsepp \emph{et
al.}, 2020; Hsu \emph{et al.}, 2021; Qie \emph{et al.}, 2024), but was
not intended to mimic natural environmental variability. Because the
realized VPD stimulus varied among treatments, we checked whether this
could confound treatment, \fgmax{} components, and \gcl{} effects on
stomatal closure kinetics using curve-level VPD covariates (Pichaco
\emph{et al.}, 2024). We found that our conclusions were robust to this
potential confound (Notes S1; Figs.
S1--S4). Following the transient
WWR, we started logging data once \gsw{} declined back to its pre-step
value (\gi) and continued logging data until \gsw{} reached its nadir
(Figure~\ref{fig-conceptual}a). We then acclimated the leaf to low light
(\ppfdequals{150}) and \rhequals{70} before inducing stomatal closure
with low \rh{} and logging data again, on average 48.9 minutes after the
start of the high-light curve.

Prior to analysis, we removed unreliable and unusable data points,
removed points not at instrument steady state, removed the hysteretic
portion of each curve at low \gsw{}, removed outliers within each curve,
removed replicates with no overlap between amphi and pseudohypo curves
from the same leaf, and thinned redundant data points within each curve.
The rationale and procedure for each cleaning step is described in Muir
\emph{et al.} (2025). Additionally, we excluded stomatal kinetic
parameter estimates from five curves where \(\tau\) was estimated to be
greater than \(\SI{1097}{\second}\), far exceeding the typical estimates
(see Results). After these steps, the total number of humidity response
curves analyzed is given in Table~S3. The median
number of curves per population per treatment was 9 and the median
number of points per curve was 26. We extracted the posterior
distribution for \(\tau\) and \(\lambda\) from each fitted curve to
calculate the parameter estimate (median) and uncertainty (standard
deviation) for subsequent analyses.

\subsubsection{\texorpdfstring{Stomatal anatomical traits \gcl{} and
\fgmax{}}{Stomatal anatomical traits  and }}\label{stomatal-anatomical-traits-and}

We estimated the stomatal density and \gcl{} on ab- and adaxial surfaces
from all leaves used for humidity response curves using standard methods
described fully in Muir \emph{et al.} (2025). We excluded three
individuals with extreme outlying guard cell length values. Preliminary
analysis of variation in guard cell length among growth treatments, leaf
types, and accessions identified these outliers based on a residual
value greater than 0.45, which visually appeared as a clear break in the
residual distribution. When the leaflet used for the LI-6800
humidity-response curve did not have a matching leaflet-level anatomical
measurement, we used the first available anatomical measurement from a
different leaflet on the same leaf as a proxy; this fallback was used
for 9 of 561 individual-by-leaf-type combinations (1.6\%). For each
surface we calculated the anatomical maximum stomatal conductance
(\gmax) at a reference leaf temperature of \(\SI{25}{\degreeCelsius}\)
following Sack \& Buckley (2016) as:

\[g_\mathrm{max} = b m d l_\mathrm{gc} / \sqrt{2}.\]

The biophysical and morphological constants \(b\) and \(m\) are:

\[b = \frac{D_\mathrm{wv}}{v},~\mathrm{and}\]\\
\[m = \frac{\pi c ^ 2}{j^{0.5} (4 h j + \pi c)},\]

where \(D_\mathrm{wv}\) is the diffusion coefficient of water vapor in
air and \(v\) is by the kinematic viscosity of dry air. We assumed
\(D_\mathrm{wv} = \SI{2.49e-5}{\meter\squared\per\second}\) and
\(v = \SI{2.24e-2}{\raiseto{3}\meter\per\mol}\) (Monteith \& Unsworth,
2013). For kidney-shaped guard cells like those in wild tomatoes,
\(c = h = j = 0.5\). To estimate \fgmax, we divided the initial \gsw{}
estimated from the parameter \gi{} in Equation~\ref{eq-weibull} from
each humidity response curve by the anatomical \gmax{} for that leaf.
The estimated \gi{} was very close to the maximum \gsw{} observed near
the beginning of each curve (Figure~S5), suggesting
that the estimated \gi{} was a good proxy for stomatal conductance at
the time of the step change in humidity. Because \gi{} is estimated from
the same nonlinear fit as \(\tau\) and \(\lambda\), we used a null
simulation to confirm that this does not, by itself, induce a spurious
\fgmax{}-\(\tau\) association (Notes S2). We used the anatomical \gmax{}
for the surface(s) measured in each curve to calculate \fgmax. For amphi
curves, we used the sum of adaxial and abaxial \gmax{}; for pseudohypo
curves, we used only the abaxial \gmax{}. Although \gmax{} calculated
this way is a theoretical upper bound based on stomatal geometry rather
than a directly measured quantity, it corresponds well to empirically
measured maximum \gsw{} in experimentally controlled \emph{Arabidopsis}
(Dow \emph{et al.}, 2014).

\subsubsection{Estimating kinetic
parameters}\label{estimating-kinetic-parameters}

We estimated \(\tau\) and \(\lambda\) for each response curve by fitting
Equation~\ref{eq-weibull} to a time course of \gsw{} values using
Bayesian nonlinear regression. All models were fit using Bayesian HMC
sampling in the probabilistic programming language \emph{Stan} (Stan
Development Team, 2026) using the \emph{R} package \textbf{brms} version
2.23.0 (Bürkner, 2017). We used \emph{CmdStan} version 2.38.0 and
\textbf{cmdstanr} version 0.9.0 (Gabry \emph{et al.}, 2025) to interface
with \emph{R} version 4.6.1 (R Core Team, 2026). We sampled 4,000
post-warmup draws from the posterior distribution across four chains. We
adjusted the number of sampling iterations and thinning interval for
each curve so that the Gelman-Rubin convergence statistic (\(\hat{R}\))
(Gelman \& Rubin, 1992) was less than 1.05 and the bulk effective sample
size (ESS) was greater than 400 for all model parameters, and there were
fewer than 10 divergent transitions. To answer Question 1, we assessed
model fit using Bayesian \(R^2\) (Gelman \emph{et al.}, 2019)
implemented in \textbf{brms}.

\subsubsection{Estimating effects of stomatal anatomy on kinetic
parameters}\label{estimating-effects-of-stomatal-anatomy-on-kinetic-parameters}

We first used model selection to determine whether anatomy (\gcl{} and
\fgmax) affected kinetic parameters (\(\lambda\) and \(\tau\); Question
2) and at what level of biological organization (Question 3). To infer
the effect of \fgmax{}, we modeled effects of component traits \loggi{}
and \loggmax{} directly because
\(\log \left( f_\mathrm{gmax} \right) = \log \left( g_\mathrm{i} \right) - \log \left( g_\mathrm{max} \right)\).
If \fgmax{} \emph{per se} affects kinetics then both \loggi{} and
\loggmax{} should be in the selected model. Alternatively, \gi{} or
\gmax{} could effect kinetic parameters independent of \fgmax. From the
selected model, we used parameter estimates and 95\% confidence
intervals to test the predictions of our hypotheses. We interpreted
nonzero fixed effects of \gcl{}, \gi{}, and \gmax{} on \(\lambda\) and
\(\tau\) as evidence that these anatomical traits influence variation in
stomatal kinetics among individuals. We interpreted nonzero phylogenetic
correlations among traits as evidence that these traits covary among
populations. We predicted that \gcl{} and \gi{} should be positively
associated with \(\lambda\) and \(\tau\), whereas the effect of \gmax{}
should be negative. Finally, we used path analysis to test whether
treatment-associated plasticity in \(\tau\) was statistically consistent
with an indirect path through \gcl{} and/or \fgmax{} components.
Table~\ref{tbl-criteria} summarizes the statistical criteria for
determining whether and at what level of biological organization
stomatal anatomical traits affected kinetic parameters. The subsections
below describe specific procedures for deciding whether those are met.

\begin{table}

\caption{\label{tbl-criteria}Criteria used to determine whether and at
what level of biological organization stomatal anatomical traits
affected kinetic parameters.}

\centering{

\centering
\begin{tabular}{>{\centering\arraybackslash}m{1.275in}>{\raggedright\arraybackslash}m{4.5in}}
\toprule
Level of organization & Criteria supporting an association between stomatal anatomy and kinetics\\
\midrule
\cellcolor{gray!10}{Among individuals} & \cellcolor{gray!10}{\begin{itemize} \item{Anatomical parameters as fixed effects improve model fit (lower LOOIC)} \item{Anatomical effects on kinetics are in the predicted direction} \item{Confidence intervals of anatomical effects do not overlap zero} \end{itemize}}\\
Plasticity among individuals & \begin{itemize} \item{Plasticity in anatomy is statistically consistent with an indirect path underlying plasticity in stomatal kinetics} \end{itemize}\\
\cellcolor{gray!10}{Among populations} & \cellcolor{gray!10}{\begin{itemize} \item{Confidence intervals of phylogenetic correlation among traits do not overlap zero} \item{Phylogenetic covariance is in the predicted direction} \end{itemize}}\\
\bottomrule
\end{tabular}

}

\end{table}%

\paragraph{Bayesian multiresponse phylogenetic
models}\label{bayesian-multiresponse-phylogenetic-models}

Multiresponse models can disentangle whether traits are associated among
populations, individuals, and/or treatments that induce plasticity. The
response variables in our models are \(\tau\), \(\lambda\), \gcl, \gi,
and \gmax. All models included a subset of explanatory variables that
were integral to the experimental design or necessary to account for
statistical nonindependence. As explanatory variables, we included fixed
effects of manipulative treatments: growth light intensity (sun and
shade); measurement light intensity (high and low); and leaf type (amphi
and pseudohypo). All models also included population as both a
phylogenetically structured random effect and a non-phylogenetic random
effect following Halliwell \emph{et al.} (2025). Prior to fitting, the
phylogenetic (co)variance matrix was rescaled as a correlation matrix
(De Villemereuil \& Nakagawa, 2014). For traits estimated from multiple
curves within the same leaf (\(\lambda\), \(\tau\), \gi) we included a
random effect of individual to account for repeated measures. We
accounted for uncertainty in \(\tau\) and \(\lambda\) using the standard
deviation of the posterior distribution of each estimate (see `Humidity
response curves' for details). We modeled residual (co)variance, a
combination of among-individual (co)variance and measurement error, as a
multivariate Student \(t\)-distribution, which is more robust to extreme
values than the multivariate normal distribution (Gelman \emph{et al.},
2014). To account for heteroskedasticity, bounded measurements, and to
linearize relationships all response variables were log-transformed
prior to analysis. All models were fit with \emph{Stan} on three chains
using the same methods and convergence criteria described above for
fitting Equation~\ref{eq-weibull} to the humidity response curves.

\paragraph{Phylogenetic correlation among
traits}\label{phylogenetic-correlation-among-traits}

We inferred that traits covary among populations if the 95\% confidence
intervals of the phylogenetic correlation between traits did not overlap
zero and the correlation was in the predicted direction. For example, a
positive phylogenetic correlation between \gcl{} and \(\tau\) would be
consistent with the prediction that species with larger guard cells have
slower stomatal kinetics. Phylogenetic covariance specifically indicates
long-term evolutionary covariation among traits because our model
simultaneously accounts for trait associations among individuals, shared
plastic responses, and residual covariance. To test for possible
causation, we used the inferred phylogenetic trait covariance matrix to
estimate partial correlations between anatomical and kinetic variables
(Halliwell \emph{et al.}, 2025). We interpret a significant partial
correlation between variables as a plausible causal hypothesis.

\paragraph{Model selection and parameter estimation of individual-level
effects}\label{model-selection-and-parameter-estimation-of-individual-level-effects}

To infer whether individual-level variation in \gcl{} and \fgmax{}
influence kinetic parameters, we used the leave-one-out cross-validation
information criterion (LOOIC) to compare the fit of models using the
\emph{R} package \textbf{loo} version 2.10.1 (Vehtari \emph{et al.},
2017). We compared models with all permutations of \gcl, \gi, and
\gmax{} influencing \(\lambda\) and \(\tau\). We considered all models
within two \(\Delta\)LOOIC standard errors of the minimum LOOIC
(Table~S4) to be in the set of plausible models to
generate posterior predictions for hypothesis testing. Within each
plausible model, we assessed whether a fixed effect was significantly
different than zero based on the 95\% confidence intervals. Among the
set of plausible models, we selected the model with the highest
proportion of significant focal parameters for further analysis. Focal
parameters are the fixed and phylogenetic effects of \gcl, \gi, and
\gmax{} on \(\lambda\) and \(\tau\). This approach aids interpretability
because it favors simpler models that exclude additional nonsignificant
parameters. Because this choice is somewhat arbitrary and not based on
rigorous statistical principles, we discuss alternative interpretations
based on the other plausible models. We report Pareto \(\hat{k}\), PSIS
effective sample size, and Monte Carlo standard error diagnostics for
the plausible models in Table~S5. Nearly every
curve had good Pareto \(\hat{k}\) (\(\le 0.7\)), but the minimum PSIS
effective sample size per curve was sometimes low (as low as 94 draws
for some models), so the PSIS-LOO estimates for a small number of
individual curves, and by extension the model-level Monte Carlo standard
error, should be interpreted with some caution.

\paragraph{Path analysis to test for plastic
effects}\label{path-analysis-to-test-for-plastic-effects}

We used path analysis (Pearl, 2022) to test whether treatment-associated
plasticity in \(\tau\) was statistically consistent with an indirect
path through plasticity in stomatal traits, in response to measurement
light intensity, growth light intensity, and curve-type treatments. We
focused on \gi{} because it was the only individual-level trait
significantly associated with \(\tau\) in the selected model (see
Results). We do not treat this as a randomized-treatment causal
mediation analysis because \gi{} is not randomized, is estimated from
the same curve as \(\tau\), and could share unmeasured confounders
(e.g., hydraulic status, measurement order, realized VPD trajectory)
with \(\tau\), violating the sequential ignorability assumption required
for a causal interpretation. We assessed the plausibility of this
assumption directly with a sensitivity analysis (Notes S3). We inferred
the direct effect of treatments on \(\tau\) from model coefficients. We
estimated the indirect effect of treatments as the product of the effect
of treatment on \gi{} and the effect of that trait on \(\tau\) (Imai
\emph{et al.}, 2010). We determined significant direct and indirect
effects based on whether the 95\% confidence intervals of the effects
were different than zero. We calculated effect sizes as the point
estimate (median of the posterior) divided by the standard error
(standard deviation of the posterior).

\subsubsection{Variance decomposition and phylogenetic
heritability}\label{variance-decomposition-and-phylogenetic-heritability}

One explanation for why traits have effects at different levels of
organization is that variance in those traits is concentrated at those
different levels of organization. To evaluate this, we decomposed the
variance in response variables (\(\lambda\), \(\tau\), \gcl, \gi, \gmax)
into phylogenetic, population (nonphylogenetic), and individual levels.
The phylogenetic variance component quantifies trait evolution between
populations congruent with the phylogenetic relationships, whereas the
population (nonphylogenetic) component quantifies variance among
populations independent of phylogenetic relationships (Lynch, 1991; De
Villemereuil \& Nakagawa, 2014). We estimated these components from the
associated random-effect standard deviations for each response variable.
The among-individual variance components quantify variance among
individuals within the same population and treatment. We estimated the
among-individual component as the sum of the individual random-effect
and residual standard deviations. It was not possible to estimate
individual random-effects in \gcl{} or \gmax{} with our study design
because we only measured average guard cell length and \gmax{} once per
leaf. The among-individual component is biased upwards because it also
includes measurement error, which cannot be separately estimated with
our experimental design. The phylogenetic heritability is equivalent to
the phylogenetic variance component and is estimated following Lynch
(1991) and De Villemereuil \& Nakagawa (2014) as:

\[h^2_{\mathrm{phy}} = \frac{\sigma^2_\mathrm{phy}}{\sigma^2_\mathrm{phy} + \sigma^2_\mathrm{pop} + \sigma^2_\mathrm{ind}},\]

where \(\sigma^2_\mathrm{phy}\), \(\sigma^2_\mathrm{pop}\),
\(\sigma^2_\mathrm{ind}\) are the phylogenetic, population
(nonphylogenetic), and among-individual variance components,
respectively. Because phylogenetic heritability was estimated after
accounting for fixed treatment effects, estimates obtained in controlled
environments likely overestimate phylogenetic heritability in natural
environments, where plasticity is a larger component of trait variation.

\subsubsection{Effect of adaxial stomata on
kinetics}\label{effect-of-adaxial-stomata-on-kinetics}

If adaxial stomata respond faster than abaxial stomata (Question 4), we
predicted that \(\tau\) should be greater in pseudohypo leaves compared
to amphi. We evaluated this prediction by estimating the fixed effect of
leaf type (amphi vs pseudohypo) on \(\tau\) in the multiresponse model.
We did not estimate adaxial-only kinetics because it was not part of the
original experimental design (Muir \emph{et al.}, 2025).

\subsection{Results}\label{results}

\subsubsection{\texorpdfstring{The time constant \(\tau\) explains most
variation in stomatal closure
kinetics}{The time constant \textbackslash tau explains most variation in stomatal closure kinetics}}\label{the-time-constant-tau-explains-most-variation-in-stomatal-closure-kinetics}

Stomatal conductance (\gsw) decreased rapidly in response to a step
change in humidity. The model in Equation~\ref{eq-weibull} fit the data
extremely well (Figure~S6). The mean and minimum
Bayesian \(R^2\) across all curves was 0.9994 and 0.9803, respectively.
Consider a typical leaf in the experiment characterized by the median
time constant (\(\tau\)) and lag time (\(\lambda\)) among wild tomato
populations in all treatments. After humidity decreased, the transient
WWR elapsed, and \gsw{} returned to the initial value \gi, it took on
average 128 s for \gsw{} to decrease half-way (\(t_{50}\)) from its
initial to final (\gf) steady state value. In terms of the kinetic model
parameters, most variation among leaves was due to difference in the
time constant rather than the lag time. In the previous example,
increasing \(\lambda\) from its median to its maximum estimated value
among populations increased the \(t_{50}\) by only
\(\SI{7.1}{\second}\). In contrast, increasing \(\tau\) from its median
to maximum value increased \(t_{50}\) by \(\SI{86}{\second}\). Hence, we
primarily focus on treatments and anatomical parameters affecting
\(\tau\).

\subsubsection{\texorpdfstring{Variation in guard cell length and
\fgmax{}
components}{Variation in guard cell length and  components}}\label{variation-in-guard-cell-length-and-components}

Guard cell length (\gcl) differed among populations and varied to a
lesser extent between growth light intensity treatments and leaf
surfaces. In contrast, stomatal conductance as a fraction of maximum
conductance (\fgmax) was most strongly determined by measurement light
intensity. Because \gmax{} is unaltered by measurement light intensity,
this treatment effect was entirely caused by changes in \gi. Excluding
treatment effects, phylogenetically structured differences among
populations (i.e., phylogenetic heritability) accounted for 61.9\%
{[}95\% CI: 18.9\% to 82.8\%{]} of the variance in \gcl{}
(Figure~\ref{fig-variance}, Table~S6). There was also
developmental plasticity in \gcl. On average, \gcl{} increased by
\(12.1\)\% {[}95\% CI: \(11.1\)\% to \(13.1\)\%{]} in the sun treatment
(Figure~\ref{fig-accession-anatomy}a; Table~S7). The
average \gcl{} of pseudohypo leaves was shorter by \(-5.8\)\% {[}95\%
CI: \(-6.6\)\% to \(-5.0\)\%{]} compared to amphi leaves
(Figure~\ref{fig-accession-anatomy}a; Table~S7). This was
not because of plasticity, as we measured amphi and pseudohypo curves on
the same leaf. Rather, this was a composition effect. Adaxial stomata
are larger than abaxial stomata in these species (Figure~S7),
so the average \gcl{} in pseudohypo leaves was lower than in amphi
leaves because the former only included abaxial stomata.

The phylogenetic heritability of \fgmax{} components was low (\gi) to
moderate (\gmax), accounting for 16.9\% {[}95\% CI: 1.1\% to 37.9\%{]}
and 59.1\% {[}95\% CI: 34.4\% to 75.5\%{]} of the variance, respectively
(Figure~\ref{fig-variance}, Table~S6). High measurement
light intensity increased \gi{} by an average of \(113.2\)\% {[}95\% CI:
\(107.7\)\% to \(118.5\)\%{]} (Figure~\ref{fig-accession-anatomy}b;
Table~S7). Sun leaves had on average \(83.3\)\% {[}95\% CI:
\(79.5\)\% to \(87.1\)\%{]} greater \gmax{} than shade-grown plants
(Figure~\ref{fig-accession-anatomy}b; Table~S7). The
pseudohypo leaf type had on average \(-19.5\)\% {[}95\% CI: \(-21.5\)\%
to \(-17.3\)\%{]} lower \gi{} than untreated, amphistomatous leaves
(Figure~\ref{fig-accession-anatomy}b; Table~S7). This is
presumably because a large fraction of stomata were blocked by the
pseudohypo treatment and increased aperture of the abaxial stomata did
not fully compensate. This observation differs from wheat, where \gsw{}
on one surface was not affected by blocking gas exchange in the other
(Wall \emph{et al.}, 2022).

\begin{figure}

\centering{

\Large{[Figure 2 goes here]}

}

\caption{\label{fig-variance}\textbf{Traits vary at different levels of biological organization}.
A large fraction of the variation in guard cell length (\gcl),
anatomical maximum stomatal conductance (\gmax), and the time constant
\(\tau\) is phylogenetically structured among \textit{Solanum}
populations. In contrast, most of the variance in initial stomatal
conductance (\gi) and the lag time (\(\lambda\)) is among individuals.
Nonphylogenetic variation among populations was low for all traits.}

\end{figure}%

\begin{figure}

\centering{

\Large{[Figure 3 goes here]}

}

\caption{\label{fig-accession-anatomy}Caption on following page.}

\end{figure}%

\addtocounter{figure}{-1}

\begin{figure}

\centering{

\mbox{}

}

\caption{\label{fig-accession-anatomy-cap}(previous page)
\textbf{Guard cell length (\gcl) and stomatal conductance as a fraction of maximum conductance (\fgmax) vary among \textit{Solanum} populations and treatments.}
In both panels, black points are mean trait value per population within
a particular treatment. Gray lines connect values within populations
across treatments. (a) Much of the variation in \gcl{} is among
populations, but developmental plasticity to growth light intensity
tends to increase \gcl{} in sun (right facet) compared to shade (left
facet) grown plants. There is no plasticity to leaf type (amphi
vs.~pseudohypo), but the average \gcl{} is slightly lower in pseudohypo
leaves because adaxial stomata tended to be larger than abaxial stomata
on the same leaf. (b) Measurement light intensity has the most direct
effect on \fgmax{} because initial stomatal conductance (\gi) increases
under high light intensity. To a lesser extent, \fgmax{} varies among
populations, is slightly lower in sun-grown plants, and slightly higher
in pseudohypo leaf types.}

\end{figure}%

\subsubsection{Variation in stomatal closure kinetic
parameters}\label{variation-in-stomatal-closure-kinetic-parameters}

Estimates of stomatal closure kinetic parameters for each curve are in
Table~S8 and estimates for each population are in
Table~S9. The time constant \(\tau\)
exhibited moderate phylogenetic heritability, accounting for 39.8\%
{[}95\% CI: 18.7\% to 60.4\%{]} of the variance, whereas \(\lambda\)
varied more within populations
(Figure~\ref{fig-variance}, Table~S6). Some light
treatments also affected kinetic parameters. The \(\tau\) in sun plants
was -10.7\% {[}95\% CI: -15.1\% to -6.06\%{]} lower than shade plants
after accounting for other explanatory variables
(Figure~\ref{fig-accession-kinetics}, Table~S7). Within the
same leaf, increasing the measurement light intensity from \ppfdqty{150}
to \ppfdqty{2000} significantly increased \(\tau\) by 8.6\% {[}95\% CI:
1.4\% to 16.4\%{]}
(Figure~\ref{fig-accession-kinetics}, Table~S7). The lag
time \(\lambda\) responded to growth, but not measurement light
intensity after accounting for other explanatory variables
(Figure~\ref{fig-accession-kinetics}, Table~S7). In sun
plants, \(\lambda\) was on average 3.43\% {[}95\% CI: 1.4\% to 5.67\%{]}
greater than in shade plants. High measurement light had little effect
on \(\lambda\), with an estimated change of -1.55\% {[}95\% CI: -4.35\%
to 1.39\%{]}.

\begin{figure}

\centering{

\Large{[Figure 4 goes here]}

}

\caption{\label{fig-accession-kinetics}\textbf{Stomatal kinetics parameters $\lambda$ (lag time; top facets) and $\tau$ (time constant; lower facets) vary among wild tomato populations.}
Stomatal conductance decreases faster (lower \(\tau\)) in response to a
step change in vapor pressure deficit (VPD) in leaves when measured
under low light intensity. The pattern was consistent for both shade
(left facets) and sun (right facets) grown plants. The pattern for
\(\lambda\) was qualitatively similar to that for \(\tau\). Each point
is the average parameter value for one accession in that treatment
combination. The growth and measurement light intensity treatments are
described in the Materials and Methods section.}

\end{figure}%

\subsubsection{\texorpdfstring{Guard cell length and \gi{} influence
stomatal
kinetics}{Guard cell length and  influence stomatal kinetics}}\label{guard-cell-length-and-influence-stomatal-kinetics}

Guard cell length and \gi{}, a component of \fgmax, affected stomatal
closure kinetics in the directions predicted by our hypotheses, but they
operated at different levels of biological organization. Among the
models we tested, six were plausible in that their predictive accuracy
was similar based on \(\Delta \mathrm{LOOIC}\)
(Table~S4). All models in the plausible set indicated
among-individual level effects of \gi{} on \(\tau\) and \(\lambda\), as
we discuss in more detail below. Models that included both individual
and phylogenetic effects of \gcl{} had similar LOOIC values but the
selected model retained only phylogenetic effects because our selection
criteria favored models that excluded nonsignificant parameters
(Table~S4). In models that included both individual
and phylogenetic effects, neither the phylogenetic correlation nor the
fixed effect regression coefficient were significantly different than
zero (Figure~S8; Figure~S9).
These parameters are highly collinear in the posterior
(Figure~S10) indicating difficulty in estimating both
terms simultaneously and ``smearing'' out their confidence intervals
(McElreath, 2020). Hence, the model with only phylogenetic effects is
easier to interpret. This model is also consistent with the high
phylogenetic heritability in both \gcl{} and \(\tau\)
(Figure~\ref{fig-variance}). Visually, we observed positive
relationships among population mean trait values
(Figure~\ref{fig-gcl-tau}) but little relationship between \gcl{} and
\(\tau\) within populations (results not shown). When comparing many
models simultaneously, differences in LOOIC performance can occur by
chance (McLatchie \& Vehtari, 2024), so judgement calls among models
with similar performance can be necessary. The preponderance of evidence
indicates that \gcl{} explains variation primarily among rather than
within our populations. The results in the main text are based on
estimates from the selected model. Individual-level fixed effects
(Figure~S8) and phylogenetic correlations
(Figure~S9) of \gcl, \gi, and \gmax{} on \(\tau\)
and \(\lambda\) were qualitatively consistent across all plausible
models.

Leaves with greater \gi{} closed more slowly than those with lower \gi{}
(Figure~\ref{fig-gi-tau}; Table~S7). For example,
increasing \gi{} from 0.2 to
\(\SI{0.4}{\mole\per\meter\squared\per\second}\) resulted in a 25.1\%
{[}95\% CI: 17.6\% to 32.6\%{]} increase in \(\tau\) for a typical leaf
in the reference treatment set (shade, low, amphi). We found little
evidence that \gmax{} affected \(\tau\) (Figure~S11).
Leaves with greater \gi{} also had a significantly greater lag time,
\(\lambda\) (Figure~S12), but there was little effect of
\gmax{} (Figure~S13).

Populations that tended to have greater \gcl{} also tended to have
greater \(\tau\) as evidenced by a positive phylogenetic correlation
(Table~S10, Figure~\ref{fig-gcl-tau}). However, we found
little evidence that variation in \gcl{} explained \(\tau\) or
\(\lambda\) within populations. There were no other significant
phylogenetic correlations (Table~S10). The only marginally
significant partial correlation was between \gcl{} and \(\tau\)
(\(\rho = 0.71\) {[}95\% CI: \(-0.17\) to \(0.98\){]}), indicating a
plausible causal effect of \gcl{} on \(\tau\) among populations.

\begin{figure}

\centering{

\Large{[Figure 5 goes here]}

}

\caption{\label{fig-gcl-tau}\textbf{Phylogenetic covariance between guard cell length (\gcl, $x$-axis) and the stomatal closure time constant ($\tau$, $y$-axis) indicates that \text{Solanum} populations with larger stomata tend to close more slowly in response to humidity.}
The overall pattern was consistent across growth light intensity
treatments (left and right facets), measurement light intensity
treatments (top and bottom facets), and leaf types (point colors). The
direction and magnitude of the ellipses are such that they should
encompass 95\% of the phylogenetically structured (co)variance in \gcl{}
and \(\tau\) on log-transformed scales. The location of ellipses are set
to the median trait values in each treatment combination for visual
comparison. Each point is the average parameter value for one population
in that treatment combination. The growth and measurement light
intensity treatments are described in the Materials and Methods
section.}

\end{figure}%

\begin{figure}

\centering{

\Large{[Figure 6 goes here]}

}

\caption{\label{fig-gi-tau}\textbf{Individual-level variation in initial stomatal conductance (\gi) influences stomatal closure kinetics.}
As \gi{} (\(x\)-axis) increases, the stomatal closure time constant
(\(\tau\), \(y\)-axis) increases. The overall pattern was consistent
across grow light intensity treatments (left and right facets),
measurement light intensity treatments (top and bottom facets), and leaf
types (point colors) in \textit{Solanum} populations. Lines are
estimated using linear regression along with 95\% confidence ribbons.
The growth and measurement light intensity treatments are described in
the Materials and Methods section.}

\end{figure}%

\subsubsection{\texorpdfstring{Plasticity in \(\tau\) is an indirect
effect of plasticity in
\gi}{Plasticity in \textbackslash tau is an indirect effect of plasticity in }}\label{plasticity-in-tau-is-an-indirect-effect-of-plasticity-in}

We imposed three treatments that influenced stomatal kinetics: growth
light intensity, measurement light intensity, and leaf type. All three
treatments directly affected \(\tau\) and had indirect effects mediated
by \gi{} (Figure~\ref{fig-mediation}). We did not estimate effects
mediated by \gcl{} because the selected model did not retain direct
effects of \gcl{} on \(\tau\) (Table~S4). The indirect
effect of greater \gi{} in sun leaves increased \(\tau\) by \(11.78\)\%
{[}95\% CI: \(8.27\)\% to \(15.99\)\%{]}, counteracting the direct
negative effect. The direct effect of the pseudohypo treatment increased
\(\tau\) by \(8.60\)\% {[}95\% CI: \(5.62\)\% to \(11.56\)\%{]}, but the
indirect effect mediated by lower \gi{} changed \(\tau\) by \(-6.72\)\%
{[}95\% CI: \(-8.64\)\% to \(-4.83\)\%{]}. In contrast, direct and
indirect effects of high measurement light intensity increased \(\tau\).
The direct effect increased \(\tau\) by \(8.56\)\% {[}95\% CI:
\(1.40\)\% to \(16.40\)\%{]} and the indirect effect mediated by \gi{}
increased \(\tau\) by \(27.62\)\% {[}95\% CI: \(19.18\)\% to
\(36.24\)\%{]}, for a combined total effect of \(38.59\)\% {[}95\% CI:
\(35.29\)\% to \(41.74\)\%{]} (note that direct and indirect percentage
effects combine additively on the log scale before back-transformation,
so they do not simply sum on the percentage scale).

\begin{figure}

\centering{

\Large{[Figure 7 goes here]}

}

\caption{\label{fig-mediation}\textbf{Treatment-associated plasticity in the time constant $\tau$ is statistically consistent with an indirect path through initial stomatal conductance (\gi) in \textit{Solanum} populations.}
The three treatments are growth light intensity (top facet; shade
vs.~sun), leaf type (middle facet; amphi vs.~pseudohypo), and
measurement light intensity (bottom facet; low vs.~high). Horizontal
arrows at the bottom of each graph indicate the direct effect of the
treatment level on \(\tau\); slanted arrows leading to or from \gi{}
indicate indirect effects of the treatment on \(\tau\) through \gi{}.
Line thickness is proportional to the effect size and line color
indicates whether the effects are negative (red) or positive (blue).
Parameter estimates of the effect and 95\% confidence intervals are
above each line. The selected model did not retain direct effects of
\gcl{} on \(\tau\) and hence this variable could not be included in this
analysis. This path analysis assumes sequential ignorability (no
unmeasured confounder of the \gi{}-\(\tau\) relationship); see Notes S3
for a sensitivity analysis of this assumption.}

\end{figure}%

\subsubsection{\texorpdfstring{Adaxial stomata close faster than abaxial
stomata, but the effect is masked by plasticity in
\gi{}}{Adaxial stomata close faster than abaxial stomata, but the effect is masked by plasticity in }}\label{adaxial-stomata-close-faster-than-abaxial-stomata-but-the-effect-is-masked-by-plasticity-in}

Although a simple comparison of \(\tau\) between paired amphi and
pseudohypo leaf types shows little net difference (results not shown),
the selected model reveals a significant direct effect of leaf type on
\(\tau\) once \gi{} is accounted for: pseudohypo leaves, which have only
abaxial stomata, closed more slowly than amphi leaves
(Figure~\ref{fig-mediation}; Table~S7). This direct effect
is consistent with the hypothesis that adaxial stomata close faster than
abaxial stomata. However, pseudohypo leaves also had lower \gi{} and
lower \gi{} is itself associated with faster closure
(Figure~\ref{fig-gi-tau}). This indirect path through \gi{} counteracts
most of the direct effect (Figure~\ref{fig-mediation}), which is why the
raw comparison of paired leaves does not reveal much difference in
kinetics between leaf types even though adaxial stomata appear to close
faster once \gi{} is controlled for.

\subsection{Discussion}\label{discussion}

The ability for stomata to adjust to fluctuating environmental
conditions likely evolved in early land plants because of selection to
optimize water use relative to carbon gain (Raven, 2002). Hydroactive
control of stomatal aperture in angiosperms increases flexibility to
modulate carbon gain and water loss depending on leaf water status
(Buckley, 2019), perhaps enabling a competitive advantage in
environments where light and evaporative demand fluctuate at high
frequency (McAdam \& Brodribb, 2012, 2015; Pichaco \emph{et al.}, 2026).
We now understand the steady-state response of stomatal conductance
(\gsw) to global increases in vapor pressure deficit (VPD) better but
there is greater uncertainty surrounding kinetic responses to VPD
(Buckley \emph{et al.}, 2023). Here we considered two nonmutually
exclusive mechanistic hypotheses for how stomatal anatomy influences
stomatal closure kinetics. First, smaller stomata may close faster
because of greater surface area to volume ratio (Drake \emph{et al.},
2013). Second, guard cells under higher turgor pressure respond more
slowly because of reduced cell wall elasticity (Franks \emph{et al.},
2012). Based on these hypotheses, we predicted that leaves with smaller
guard cells and operating further from their anatomical maximum \gsw{}
(i.e., lower \fgmax) would have faster stomatal closure kinetics (lower
\(\tau\)). The parameter \fgmax{} is here defined as the initial
stomatal conductance (\gi) divided by \gmax. Analysis of 2,125 time
courses of stomatal closure from 29 populations in multiple environments
demonstrates that both guard cell size and \gi{}, a component of \fgmax,
influence kinetics at different levels of biological organization. This
new understanding suggests unexplored explanations for the relationships
between stomatal size, density, speed, and operational \gsw{} (\gop)
among angiosperms.

The kinetics of stomatal closure in response to VPD closely match
response to light, suggesting similar underlying mechanisms. The
lag-time model proposed by Woning \emph{et al.} (2026) for light
responses also describes humidity responses extremely well following the
passive wrong-way response (WWR). The fitted models explained on average
99.94\% of the variation in \gsw{} during closure
(Figure~S6). The formal mathematical similarity of
light and humidity response might indicate that certain factors, such as
guard cell size, influence stomatal kinetics responses in general.
Alternatively, disparate underlying processes may be described by the
same macroscale formalism. For example, humidity responses may be
controlled by ABA accumulation or loss (McAdam \emph{et al.}, 2016),
whereas \(C_\mathrm{i}\)-independent red-light responses can be
influenced by accumulation or loss of photosynthetic products (Matthews
\emph{et al.}, 2020). These disparate processes may both have similar
kinetic properties, but would not necessarily imply that VPD and light
responses are strongly correlated. Future research comparing the
correlation between light, VPD, CO\textsubscript{2}, and other stomatal
kinetic responses can determine the extent to which generic or specific
factors determine stomatal opening and closing rates (Creese \emph{et
al.}, 2014; Kübarsepp \emph{et al.}, 2020).

The positive association between \gcl{} and \(\tau\)
(Figure~\ref{fig-gcl-tau}) is consistent with the hypothesis that
smaller cells respond faster because they have a greater surface area to
volume ratio for osmolyte flux (Franks \& Farquhar, 2007; Drake \emph{et
al.}, 2013; Lawson \& Vialet-Chabrand, 2019; Woning \emph{et al.},
2026). Among several of the same \emph{Solanum} species, larger guard
cells are also associated with slower opening kinetics in response to
light (Yoshiyama \emph{et al.}, 2024), suggesting a general mechanism.
Specifically, this association implies that the maximum rate of stomatal
closure is limited by membrane surface area. We do not have the data to
evaluate the hypothesis mechanistically, but this should be tested in
the future. In addition to guard cell size, it will be useful to test
whether the ratio of guard cell to adjacent epidermal size affects the
strength of the WWR, which may be necessary to induce more rapid
stomatal closure (Buckley \emph{et al.}, 2003; Pichaco \emph{et al.},
2025). An alternative explanation is that the correlation between \gcl{}
and \(\tau\) is not causal. The fact that \gcl{} did not explain much of
the variation in \(\tau\) among individuals within populations suggests
that \gcl{} may not directly affect \(\tau\). However, the high
phylogenetic heritability and, hence, low variability in \gcl{} within
populations may also explain why we were unable to detect a
within-population effect of \gcl{} on \(\tau\).

A potential consequence of greater stomatal aperture is that the pore
closes more slowly when the environment changes. Much of the variation
in \(\tau\) among individuals within populations was explained by
variation in \gi{} (Figure~\ref{fig-gi-tau}). This association is not
simply an artifact of the magnitude of the difference between \gi{} and
\gf{}, because \(\tau\) captures the proportional rate of change, not
the absolute change, in \gsw. This factor was positively associated with
\(\tau\) among individuals within light treatments, and
treatment-associated plasticity in \(\tau\) was statistically consistent
with an indirect path through \gi{} (Figure~\ref{fig-mediation}). The
rate of stomatal opening in response to light is associated with \gop{}
in a sample of rainforest trees (Kardiman \& Ræbild, 2018), which may
indicate this is a more general phenomenon.

The association between \gi{} and \(\tau\) could be explained by the
hypothesis that guard cell wall elasticity declines at greater turgor
pressure (Franks \emph{et al.}, 2012), like a rubber band that is
stretched under tension. However, if this were the case, we would also
expect a negative association between \gmax{} and \(\tau\) because
leaves with greater \gmax{} could have the same \gi{} with a smaller
stomatal aperture. However, this effect was weakly supported
(Figure~S11, Figure~S13). The selected
model did not include an individual-level effect of \gmax{} on kinetic
parameters. Although the second-best model included \gmax{} and the
effect was in the predicted direction (Table~S4), the
effect was not significantly different than 0
(Figure~S8). The weak effect of \gmax{} compared to
\gi{} casts doubt on whether a nonlinear relationship between turgor
pressure and pore aperture is the best explanation. An alternative,
albeit not mutually exclusive explanation, is a nonlinear relationship
between pore aperture and \gsw{} that weakens the association between
these variables at high \gsw. Airspace resistance within the substomatal
cavity and stomatal interference are plausible mechanisms that would
have this effect (Ting \& Loomis, 1965; Kaiser, 2009; Lehmann \& Or,
2015), although it is not clear that they are sufficient to produce the
observed relationship between \gi{} and \(\tau\) in wild tomatoes. Our
study is limited by the fact that we did not directly measure guard cell
turgor or stomatal aperture, but future studies should use newly
developed methods (Brodersen \emph{et al.}, 2025) to directly evaluate
the association between these traits.

The impact of \gi{} was most pronounced when comparing measurement light
intensity treatments. Leaves acclimated to high light intensity had on
average 113\% greater \gi{} than the same leaves at low light intensity
in the reference treatment (shade, amphi). As a consequence, stomata
closed more slowly in leaves acclimated to high light intensity, with a
27.6\% greater time constant on average in the reference treatment,
based on parameter estimates from the selected model
(Table~S7). We did not assess whether \gi{} also affects
stomatal opening and this should be measured to further test this
hypothesis. The implication is that natural selection should favor
\gop{} values that not only optimize instantaneous carbon gain and water
loss in the present environment, but also carbon gain and water loss
integrated through time in a fluctuating environment. We predict that
rapid, temporally-uncorrelated fluctuating environments will favor
leaves with greater \gop{} to slow kinetic responses and prevent
overreacting to transitory changes that are likely to revert; slower,
temporally-correlated changes will favor lower \gop{} to enable close
tracking of the environment.

The comparison between measurement light intensity treatments should be
interpreted cautiously, however, because it was confounded with
measurement order: high light curves always preceded low light curves on
the same leaf, following a prior low-humidity closure event (see
Materials and Methods). We cannot rule out that some of the apparent
effect of measurement light intensity on \gi{} and \(\tau\), and the
corresponding indirect path through \gi{} in our mediation analysis
(Figure~\ref{fig-mediation}), reflects order, incomplete recovery, or
prior VPD exposure rather than measurement light intensity itself. This
concern applies specifically to the measurement light comparison and not
to growth light intensity or leaf type, which were randomized or
alternated, respectively.

The direct effect of leaf type on \(\tau\) indicates that adaxial
stomata close faster than abaxial stomata after accounting for
differences in \gi{}. This result is consistent with the broader
association between greater adaxial stomatal density and faster stomatal
closure across land plants (Haworth \emph{et al.}, 2018; Woning \emph{et
al.}, 2026). One plausible reason adaxial stomata may have evolved to
close intrinsically faster is because adaxial stomata directly overlie
the photosynthetically active palisade mesophyll and failure to close
quickly could rapidly desiccate this tissue and impair photosynthesis
(Tezara \emph{et al.}, 1999; Buckley \emph{et al.}, 2017a).
Hypostomatous leaves may experience less dehydration risk because the
palisade is hydraulically buffered from changing VPD when stomata are
confined to the abaxial surface (Buckley \emph{et al.}, 2017a). If
traits such as guard cell size that influence kinetics act similarly on
both surfaces, this shared basis, together with the intrinsic speed
difference between surfaces documented here, could jointly explain why
leaves and species with a greater proportion of adaxial stomata tend to
have faster overall stomatal kinetics.

Future research is needed to determine how stomatal kinetic responses
measured in controlled conditions impact performance in nature. Our
experimental protocol induced stomatal closure by abruptly switching
incoming air to near-zero humidity, a perturbation far more severe than
the gradual or partial VPD fluctuations leaves experience in the field.
This approach is well suited for estimating the maximum rate of stomatal
closure and comparing it among individuals and treatments, but it need
not follow that the same anatomical and physiological factors that
govern \(\tau\) in our experiment also govern kinetic variation under
more naturalistic conditions, where closure VPD may fluctuate more
gradually. Future studies should also control the VPD stimulus
throughout the entire response curve, rather than allowing it to vary as
transpiration itself alters chamber humidity, as it did here
(Table~S2). It will also be important to consider whether the
substomatal airspace remains saturated during the humidity step, because
\gsw{} will be underestimated if the airspace is subsaturated (Cernusak
\emph{et al.}, 2018; Buckley \& Sack, 2019; Wong \emph{et al.}, 2022;
Rockwell \emph{et al.}, 2022). We hypothesize that ecological conditions
which favor close tracking of fluctuating environments will select for
leaves with traits like smaller guard cell size and low \gi{} that
enable rapid kinetics.

\subsection{Acknowledgements}\label{acknowledgements}

Comments from two anonymous reviewers greatly improved the manuscript.
Sam McKlin and Tom Buckley helped with protocol development. Justin
Alter, Max Gatlin, Joana Kim, Jenna Matsuyama, Brandon Najarian, Dachuan
Wang, and Kai Yasuda contributed to data collection. We used Copilot,
ChatGPT, and Posit Assistant (Claude) to assist with coding and
proofreading text.

\subsection{Competing interests}\label{competing-interests}

None declared.

\subsection{Funding}\label{funding}

US National Science Foundation OIA-1929167 (C.D.M.)

\subsection{Author contributions}\label{author-contributions}

Conceptualization: C.D.M.; Data acquisition: C.D.M., W.S.L.; Data
analysis and visualization: C.D.M.; Writing -- original draft: C.D.M.;
Writing -- review \& editing: C.D.M., W.S.L.

\subsection{Data availability}\label{data-availability}

All data and code are available on GitHub
(\url{https://github.com/cdmuir/solanum-kinetics}). Data are archived on
Dryad (\url{https://doi.org/10.5061/dryad.f7m0cfzcv}); code is archived
on Zenodo (\url{https://doi.org/10.5281/zenodo.22802259}).

\subsection{References}\label{references}

\singlespacing

\phantomsection\label{refs}
\begin{CSLReferences}{0}{1}
\bibitem[\citeproctext]{ref-aasamaa_stomatal_2011}
\textbf{\textbf{Aasamaa K}, \textbf{Sõber A}}. \textbf{2011a}.
\href{https://doi.org/10.1016/j.envexpbot.2010.10.013}{Stomatal
sensitivities to changes in leaf water potential, air humidity,
{CO}{\textsubscript{2}} concentration and light intensity, and the
effect of abscisic acid on the sensitivities in six temperate deciduous
tree species}. \emph{Environmental and Experimental Botany} \textbf{71}:
72--78.

\bibitem[\citeproctext]{ref-aasamaa_responses_2011}
\textbf{\textbf{Aasamaa K}, \textbf{Sõber A}}. \textbf{2011b}.
\href{https://doi.org/10.1093/treephys/tpr078}{Responses of stomatal
conductance to simultaneous changes in two environmental factors}.
\emph{Tree Physiology} \textbf{31}: 855--864.

\bibitem[\citeproctext]{ref-aasamaa_leaf_2001}
\textbf{\textbf{Aasamaa K}, \textbf{Sõber A}, \textbf{Rahi M}}.
\textbf{2001}. \href{https://doi.org/10.1071/PP00157}{Leaf anatomical
characteristics associated with shoot hydraulic conductance, stomatal
conductance and stomatal sensitivity to changes of leaf water status in
temperate deciduous trees}. \emph{Australian Journal of Plant
Physiology} \textbf{28}: 765--774.

\bibitem[\citeproctext]{ref-bauer_stomatal_2013}
\textbf{\textbf{Bauer H}, \textbf{Ache P}, \textbf{Lautner S},
\textbf{Fromm J}, \textbf{Hartung W}, \textbf{Al-Rasheid KAS},
\textbf{Sonnewald S}, \textbf{Sonnewald U}, \textbf{Kneitz S},
\textbf{Lachmann N}, \emph{et al.}} \textbf{2013}.
\href{https://doi.org/10.1016/j.cub.2012.11.022}{The {Stomatal Response}
to {Reduced Relative Humidity Requires Guard Cell-Autonomous ABA
Synthesis}}. \emph{Current Biology} \textbf{23}: 53--57.

\bibitem[\citeproctext]{ref-brench_was_2026}
\textbf{\textbf{Brench RA}, \textbf{Wilson MJ}, \textbf{Thorne SJ},
\textbf{Fleming AJ}, \textbf{Gray JE}}. \textbf{2026}.
\href{https://doi.org/10.1111/nph.70830}{Was the evolution of faster
stomata driven by increased gas exchange rates rather than increasing
water use efficiency?} \emph{New Phytologist} \textbf{249}: 2355--2371.

\bibitem[\citeproctext]{ref-brodersen_situ_2025}
\textbf{\textbf{Brodersen CR}, \textbf{Brodribb TJ}, \textbf{Hochberg
U}, \textbf{Holbrook NM}, \textbf{McAdam SAM}, \textbf{Zailaa J},
\textbf{Huggett BA}, \textbf{Marmottant P}}. \textbf{2025}.
\href{https://doi.org/10.1073/pnas.2419887122}{In situ cavitation bubble
manometry reveals a lack of light-activated guard cell turgor modulation
in bryophytes}. \emph{Proceedings of the National Academy of Sciences}
\textbf{122}: e2419887122.

\bibitem[\citeproctext]{ref-bucher_inter_2016}
\textbf{\textbf{Bucher SF}, \textbf{Auerswald K}, \textbf{Tautenhahn S},
\textbf{Geiger A}, \textbf{Otto J}, \textbf{Müller A}, \textbf{Römermann
C}}. \textbf{2016}.
\href{https://doi.org/10.1007/s11258-016-0564-2}{Inter- and
intraspecific variation in stomatal pore area index along elevational
gradients and its relation to leaf functional traits}. \emph{Plant
Ecology} \textbf{217}: 229--240.

\bibitem[\citeproctext]{ref-buckley_how_2019}
\textbf{\textbf{Buckley TN}}. \textbf{2019}.
\href{https://doi.org/10.1111/nph.15899}{How do stomata respond to water
status?} \emph{New Phytologist} \textbf{224}: 21--36.

\bibitem[\citeproctext]{ref-buckley_kinetic_2023}
\textbf{\textbf{Buckley TN}, \textbf{Frehner EH}, \textbf{Bailey BN}}.
\textbf{2023}. \href{https://doi.org/10.1111/nph.18739}{Kinetic factors
of physiology and the dynamic light environment influence the economic
landscape of short-term hydraulic risk}. \emph{New Phytologist}
\textbf{238}: 529--548.

\bibitem[\citeproctext]{ref-buckley_sites_2017}
\textbf{\textbf{Buckley TN}, \textbf{John GP}, \textbf{Scoffoni C},
\textbf{Sack L}}. \textbf{2017a}.
\href{https://doi.org/10.1104/pp.16.01605}{The sites of evaporation
within leaves}. \emph{Plant Physiology} \textbf{173}: 1763--1782.

\bibitem[\citeproctext]{ref-buckley_hydromechanical_2003}
\textbf{\textbf{Buckley TN}, \textbf{Mott KA}, \textbf{Farquhar GD}}.
\textbf{2003}. \href{https://doi.org/10.1046/j.1365-3040.2003.01094.x}{A
hydromechanical and biochemical model of stomatal conductance}.
\emph{Plant, Cell \& Environment} \textbf{26}: 1767--1785.

\bibitem[\citeproctext]{ref-buckley_humidity_2019}
\textbf{\textbf{Buckley TN}, \textbf{Sack L}}. \textbf{2019}.
\href{https://doi.org/10.1002/ajb2.1282}{The humidity inside leaves and
why you should care: Implications of unsaturation of leaf intercellular
airspaces}. \emph{American Journal of Botany} \textbf{106}: 618--621.

\bibitem[\citeproctext]{ref-buckley_optimal_2017}
\textbf{\textbf{Buckley TN}, \textbf{Sack L}, \textbf{Farquhar GD}}.
\textbf{2017b}. \href{https://doi.org/10.1111/pce.12823}{Optimal plant
water economy: {Optimal} plant water economy}. \emph{Plant, Cell \&
Environment} \textbf{40}: 881--896.

\bibitem[\citeproctext]{ref-buckley_role_2011}
\textbf{\textbf{Buckley TN}, \textbf{Sack L}, \textbf{Gilbert ME}}.
\textbf{2011}. \href{https://doi.org/10.1104/pp.111.175638}{The {Role}
of {Bundle Sheath Extensions} and {Life Form} in {Stomatal Responses} to
{Leaf Water Status}}. \emph{Plant Physiology} \textbf{156}: 962--973.

\bibitem[\citeproctext]{ref-burkner_brms_2017}
\textbf{\textbf{Bürkner P-C}}. \textbf{2017}.
\href{https://doi.org/10.18637/jss.v080.i01}{{\textbf{brms}}: {An}
{\emph{R}} package for {Bayesian} multilevel models using
{\emph{Stan}}}. \emph{Journal of Statistical Software} \textbf{80}.

\bibitem[\citeproctext]{ref-cernusak_unsaturation_2018}
\textbf{\textbf{Cernusak LA}, \textbf{Ubierna N}, \textbf{Jenkins MW},
\textbf{Garrity SR}, \textbf{Rahn T}, \textbf{Powers HH}, \textbf{Hanson
DT}, \textbf{Sevanto S}, \textbf{Wong SC}, \textbf{McDowell NG},
\emph{et al.}} \textbf{2018}.
\href{https://doi.org/10.1038/s41598-018-25838-2}{Unsaturation of vapour
pressure inside leaves of two conifer species}. \emph{Scientific
Reports} \textbf{8}: 7667.

\bibitem[\citeproctext]{ref-cowan_oscillations_1972}
\textbf{\textbf{Cowan IR}}. \textbf{1972}.
\href{https://doi.org/10.1007/BF00388098}{Oscillations in stomatal
conductance and plant functioning associated with stomatal conductance:
{Observations} and a model}. \emph{Planta} \textbf{106}: 185--219.

\bibitem[\citeproctext]{ref-cowan_stomatal_1977}
\textbf{\textbf{Cowan IR}, \textbf{Farquhar GD}}. \textbf{1977}.
Stomatal function in relation to leaf metabolism and environment.
\emph{Symposia of the Society for Experimental Biology} \textbf{31}:
471--505.

\bibitem[\citeproctext]{ref-creese_are_2014}
\textbf{\textbf{Creese C}, \textbf{Oberbauer S}, \textbf{Rundel P},
\textbf{Sack L}}. \textbf{2014}.
\href{https://doi.org/10.1111/nph.12922}{Are fern stomatal responses to
different stimuli coordinated? {Testing} responses to light, vapor
pressure deficit, and {CO}{\textsubscript{2}} for diverse species grown
under contrasting irradiances}. \emph{New Phytologist} \textbf{204}:
92--104.

\bibitem[\citeproctext]{ref-damour_overview_2010}
\textbf{\textbf{Damour G}, \textbf{Simonneau T}, \textbf{Cochard H},
\textbf{Urban L}}. \textbf{2010}.
\href{https://doi.org/10.1111/j.1365-3040.2010.02181.x}{An overview of
models of stomatal conductance at the leaf level: {Models} of stomatal
conductance}. \emph{Plant, Cell \& Environment} \textbf{33}: 1419--1438.

\bibitem[\citeproctext]{ref-darwin_observations_1898}
\textbf{\textbf{Darwin F}}. \textbf{1898}.
\href{https://doi.org/10.1098/rspl.1898.0053}{Observations on stomata}.
\emph{Proceedings of the Royal Society of London} \textbf{63}: 413--417.

\bibitem[\citeproctext]{ref-devillemereuil_general_2014}
\textbf{\textbf{De Villemereuil P}, \textbf{Nakagawa S}}. \textbf{2014}.
\href{https://doi.org/10.1007/978-3-662-43550-2_11}{General
{Quantitative Genetic Methods} for {Comparative Biology}}. In:
Garamszegi LZ, ed. Modern {Phylogenetic Comparative Methods} and {Their
Application} in {Evolutionary Biology}. Berlin, Heidelberg: Springer
Berlin Heidelberg, 287--303.

\bibitem[\citeproctext]{ref-desai_abscisic_2025}
\textbf{\textbf{Desai S}, \textbf{Stroock A}}. \textbf{2025}.
\href{https://doi.org/10.64898/2025.12.26.696581}{Abscisic acid-mediated
water stress regulation can mechanistically explain oscillations and
water stress memory in stomatal conductance}.

\bibitem[\citeproctext]{ref-dow_integrated_2014}
\textbf{\textbf{Dow GJ}, \textbf{Bergmann DC}, \textbf{Berry JA}}.
\textbf{2014}. An integrated model of stomatal development and leaf
physiology. \emph{New Phytologist} \textbf{201}: 1218--1226.

\bibitem[\citeproctext]{ref-drake_smaller_2013}
\textbf{\textbf{Drake PL}, \textbf{Froend RH}, \textbf{Franks PJ}}.
\textbf{2013}. \href{https://doi.org/10.1093/jxb/ers347}{Smaller, faster
stomata: Scaling of stomatal size, rate of response, and stomatal
conductance}. \emph{Journal of Experimental Botany} \textbf{64}:
495--505.

\bibitem[\citeproctext]{ref-elliott-kingston_does_2016}
\textbf{\textbf{Elliott-Kingston C}, \textbf{Haworth M},
\textbf{Yearsley JM}, \textbf{Batke SP}, \textbf{Lawson T},
\textbf{McElwain JC}}. \textbf{2016}.
\href{https://doi.org/10.3389/fpls.2016.01253}{Does {Size Matter}?
{Atmospheric CO}{\textsubscript{2}} {May Be} a {Stronger Driver} of
{Stomatal Closing Rate Than Stomatal Size} in {Taxa That Diversified}
under {Low CO}{\textsubscript{2}}}. \emph{Frontiers in Plant Science}
\textbf{7}.

\bibitem[\citeproctext]{ref-franks_guard_2001}
\textbf{\textbf{Franks PJ}, \textbf{Buckley TN}, \textbf{Shope JC},
\textbf{Mott KA}}. \textbf{2001}.
\href{https://doi.org/10.1104/pp.125.4.1577}{Guard cell volume and
pressure measured concurrently by confocal microscopy and the cell
pressure probe}. \emph{Plant Physiology} \textbf{125}: 1577--1584.

\bibitem[\citeproctext]{ref-franks_effect_2001}
\textbf{\textbf{Franks PJ}, \textbf{Farquhar GD}}. \textbf{2001}.
\href{https://doi.org/10.1104/pp.125.2.935}{The {Effect} of {Exogenous
Abscisic Acid} on {Stomatal Development}, {Stomatal Mechanics}, and
{Leaf Gas Exchange} in {\emph{Tradescantia}}{ \emph{Virginiana}}}.
\emph{Plant Physiology} \textbf{125}: 935--942.

\bibitem[\citeproctext]{ref-franks_mechanical_2007}
\textbf{\textbf{Franks PJ}, \textbf{Farquhar GD}}. \textbf{2007}.
\href{https://doi.org/10.1104/pp.106.089367}{The {Mechanical Diversity}
of {Stomata} and {Its Significance} in {Gas-Exchange Control}}.
\emph{Plant Physiology} \textbf{143}: 78--87.

\bibitem[\citeproctext]{ref-franks_physiological_2012}
\textbf{\textbf{Franks PJ}, \textbf{Leitch IJ}, \textbf{Ruszala EM},
\textbf{Hetherington AM}, \textbf{Beerling DJ}}. \textbf{2012}.
\href{https://doi.org/10.1098/rstb.2011.0270}{Physiological framework
for adaptation of stomata to {CO} {\textsubscript{2}} from glacial to
future concentrations}. \emph{Philosophical Transactions of the Royal
Society B: Biological Sciences} \textbf{367}: 537--546.

\bibitem[\citeproctext]{ref-gabry_cmdstanr_2025}
\textbf{\textbf{Gabry J}, \textbf{Češnovar R}, \textbf{Johnson A},
\textbf{Bronder S}}. \textbf{2025}. \emph{Cmdstanr: {R Interface} to
{`{CmdStan}'}}.

\bibitem[\citeproctext]{ref-gelman_bayesian_2014}
\textbf{\textbf{Gelman A}, \textbf{Carlin JB}, \textbf{Stern HS},
\textbf{Dunson DB}, \textbf{Vehtari A}, \textbf{Rubin DB}}.
\textbf{2014}. \emph{Bayesian data analysis}. Boca Raton London New
York: CRC Press, Taylor \& Francis Group.

\bibitem[\citeproctext]{ref-gelman_rsquared_2019}
\textbf{\textbf{Gelman A}, \textbf{Goodrich B}, \textbf{Gabry J},
\textbf{Vehtari A}}. \textbf{2019}.
\href{https://doi.org/10.1080/00031305.2018.1549100}{R-squared for
{Bayesian Regression Models}}. \emph{The American Statistician}
\textbf{73}: 307--309.

\bibitem[\citeproctext]{ref-gelman_inference_1992}
\textbf{\textbf{Gelman A}, \textbf{Rubin DB}}. \textbf{1992}. Inference
from iterative simulation using multiple sequences. \emph{Statistical
Science} \textbf{7}: 457--472.

\bibitem[\citeproctext]{ref-givnish_sizes_1976}
\textbf{\textbf{Givnish TJ}, \textbf{Vermeij GJ}}. \textbf{1976}. Sizes
and shapes of liane leaves. \emph{American Naturalist} \textbf{110}:
743--778.

\bibitem[\citeproctext]{ref-grantz_stomatal_1986}
\textbf{\textbf{Grantz DA}, \textbf{Zeiger E}}. \textbf{1986}.
\href{https://doi.org/10.1104/pp.81.3.865}{Stomatal {Responses} to
{Light} and {Leaf-Air Water Vapor Pressure Difference Show Similar
Kinetics} in {Sugarcane} and {Soybean}}. \emph{Plant Physiology}
\textbf{81}: 865--868.

\bibitem[\citeproctext]{ref-grossiord_plant_2020}
\textbf{\textbf{Grossiord C}, \textbf{Buckley TN}, \textbf{Cernusak LA},
\textbf{Novick KA}, \textbf{Poulter B}, \textbf{Siegwolf RTW},
\textbf{Sperry JS}, \textbf{McDowell NG}}. \textbf{2020}.
\href{https://doi.org/10.1111/nph.16485}{Plant responses to rising vapor
pressure deficit}. \emph{New Phytologist} \textbf{226}: 1550--1566.

\bibitem[\citeproctext]{ref-halliwell_multiresponse_2025}
\textbf{\textbf{Halliwell B}, \textbf{Holland BR}, \textbf{Yates LA}}.
\textbf{2025}. \href{https://doi.org/10.1111/brv.70001}{Multi-response
phylogenetic mixed models: Concepts and application}. \emph{Biological
Reviews} \textbf{100}: 1294--1316.

\bibitem[\citeproctext]{ref-haworth_functional_2023}
\textbf{\textbf{Haworth M}, \textbf{Marino G}, \textbf{Materassi A},
\textbf{Raschi A}, \textbf{Scutt CP}, \textbf{Centritto M}}.
\textbf{2023}.
\href{https://doi.org/10.1016/j.scitotenv.2022.160908}{The functional
significance of the stomatal size to density relationship: {Interaction}
with atmospheric {[}{CO2}{]} and role in plant physiological behaviour}.
\emph{Science of The Total Environment} \textbf{863}: 160908.

\bibitem[\citeproctext]{ref-haworth_allocation_2018}
\textbf{\textbf{Haworth M}, \textbf{Scutt CP}, \textbf{Douthe C},
\textbf{Marino G}, \textbf{Gomes MTG}, \textbf{Loreto F}, \textbf{Flexas
J}, \textbf{Centritto M}}. \textbf{2018}.
\href{https://doi.org/10.1016/j.envexpbot.2018.04.010}{Allocation of the
epidermis to stomata relates to stomatal physiological control:
{Stomatal} factors involved in the evolutionary diversification of the
angiosperms and development of amphistomaty}. \emph{Environmental and
Experimental Botany} \textbf{151}: 55--63.

\bibitem[\citeproctext]{ref-hetherington_role_2003}
\textbf{\textbf{Hetherington AM}, \textbf{Woodward FI}}. \textbf{2003}.
\href{https://doi.org/10.1038/nature01843}{The role of stomata in
sensing and driving environmental change}. \emph{Nature} \textbf{424}:
901--908.

\bibitem[\citeproctext]{ref-homann_number_2002}
\textbf{\textbf{Homann U}, \textbf{Thiel G}}. \textbf{2002}.
\href{https://doi.org/10.1073/pnas.152324399}{The number of
{K}{\textsuperscript{+}} channels in the plasma membrane of guard cell
protoplasts changes in parallel with the surface area}.
\emph{Proceedings of the National Academy of Sciences} \textbf{99}:
10215--10220.

\bibitem[\citeproctext]{ref-hsu_raflike_2021}
\textbf{\textbf{Hsu P-K}, \textbf{Takahashi Y}, \textbf{Merilo E},
\textbf{Costa A}, \textbf{Zhang L}, \textbf{Kernig K}, \textbf{Lee KH},
\textbf{Schroeder JI}}. \textbf{2021}.
\href{https://doi.org/10.1073/pnas.2107280118}{Raf-like kinases and
receptor-like (pseudo)kinase {GHR1} are required for stomatal vapor
pressure difference response}. \emph{Proceedings of the National Academy
of Sciences} \textbf{118}: e2107280118.

\bibitem[\citeproctext]{ref-imai_general_2010}
\textbf{\textbf{Imai K}, \textbf{Keele L}, \textbf{Tingley D}}.
\textbf{2010}. \href{https://doi.org/10.1037/a0020761}{A general
approach to causal mediation analysis.} \emph{Psychological Methods}
\textbf{15}: 309--334.

\bibitem[\citeproctext]{ref-ishida_roles_1992}
\textbf{\textbf{Ishida A}, \textbf{Yamamura Y}, \textbf{Hori Y}}.
\textbf{1992}. \href{https://doi.org/10.1007/BF02347090}{Roles of leaf
water potential and soil-to-leaf hydraulic conductance in water use by
understorey woody plants}. \emph{Ecological Research} \textbf{7}:
213--223.

\bibitem[\citeproctext]{ref-jalakas_molecular_2021}
\textbf{\textbf{Jalakas P}, \textbf{Takahashi Y}, \textbf{Waadt R},
\textbf{Schroeder JI}, \textbf{Merilo E}}. \textbf{2021}.
\href{https://doi.org/10.1111/nph.17592}{Molecular mechanisms of
stomatal closure in response to rising vapour pressure deficit}.
\emph{New Phytologist} \textbf{232}: 468--475.

\bibitem[\citeproctext]{ref-jezek_guard_2021}
\textbf{\textbf{Jezek M}, \textbf{Silva-Alvim FAL}, \textbf{Hills A},
\textbf{Donald N}, \textbf{Ishka MR}, \textbf{Shadbolt J}, \textbf{He
B}, \textbf{Lawson T}, \textbf{Harper JF}, \textbf{Wang Y}, \emph{et
al.}} \textbf{2021}.
\href{https://doi.org/10.1038/s41477-021-00966-2}{Guard cell
endomembrane {Ca2}+-{ATPases} underpin a {`carbon memory'} of
photosynthetic assimilation that impacts on water-use efficiency}.
\emph{Nature Plants} \textbf{7}: 1301--1313.

\bibitem[\citeproctext]{ref-kaiser_relation_2009}
\textbf{\textbf{Kaiser H}}. \textbf{2009}.
\href{https://doi.org/10.1111/j.1365-3040.2009.01990.x}{The relation
between stomatal aperture and gas exchange under consideration of pore
geometry and diffusional resistance in the mesophyll}. \emph{Plant, Cell
\& Environment} \textbf{32}: 1091--1098.

\bibitem[\citeproctext]{ref-kardiman_relationship_2018}
\textbf{\textbf{Kardiman R}, \textbf{Ræbild A}}. \textbf{2018}.
\href{https://doi.org/10.1093/treephys/tpx149}{Relationship between
stomatal density, size and speed of opening in {Sumatran} rainforest
species}. \emph{Tree Physiology} \textbf{38}: 696--705.

\bibitem[\citeproctext]{ref-kerr_harking_1998}
\textbf{\textbf{Kerr NL}}. \textbf{1998}.
\href{https://doi.org/10.1207/s15327957pspr0203_4}{{HARKing}:
{Hypothesizing After} the {Results} are {Known}}. \emph{Personality and
Social Psychology Review} \textbf{2}: 196--217.

\bibitem[\citeproctext]{ref-knapp_stomatal_1990}
\textbf{\textbf{Knapp AK}, \textbf{Smith WK}}. \textbf{1990}.
\href{https://doi.org/10.1111/j.1399-3054.1990.tb08731.x}{Stomatal and
photosynthetic responses to variable sunlight}. \emph{Physiologia
Plantarum} \textbf{78}: 160--165.

\bibitem[\citeproctext]{ref-kubarsepp_are_2020}
\textbf{\textbf{Kübarsepp L}, \textbf{Laanisto L}, \textbf{Niinemets Ü},
\textbf{Talts E}, \textbf{Tosens T}}. \textbf{2020}.
\href{https://doi.org/10.1111/nph.16159}{Are stomata in ferns and allies
sluggish? {Stomatal} responses to {CO}{\textsubscript{2}}, humidity and
light and their scaling with size and density}. \emph{New Phytologist}
\textbf{225}: 183--195.

\bibitem[\citeproctext]{ref-lammertsma_global_2011}
\textbf{\textbf{Lammertsma EI}, \textbf{Boer HJD}, \textbf{Dekker SC},
\textbf{Dilcher DL}, \textbf{Lotter AF}, \textbf{Wagner-Cremer F}}.
\textbf{2011}. \href{https://doi.org/10.1073/pnas.1100371108}{Global
{CO}{\textsubscript{2}} rise leads to reduced maximum stomatal
conductance in {Florida} vegetation}. \emph{Proceedings of the National
Academy of Sciences} \textbf{108}: 4035--4040.

\bibitem[\citeproctext]{ref-lange_responses_1971}
\textbf{\textbf{Lange OL}, \textbf{Lösch R}, \textbf{Schulze E-D},
\textbf{Kappen L}}. \textbf{1971}.
\href{https://doi.org/10.1007/BF00386887}{Responses of stomata to
changes in humidity}. \emph{Planta} \textbf{100}: 76--86.

\bibitem[\citeproctext]{ref-lawson_stomatal_2014}
\textbf{\textbf{Lawson T}, \textbf{Blatt MR}}. \textbf{2014}.
\href{https://doi.org/10.1104/pp.114.237107}{Stomatal {Size}, {Speed},
and {Responsiveness Impact} on {Photosynthesis} and {Water Use
Efficiency}}. \emph{Plant Physiology} \textbf{164}: 1556--1570.

\bibitem[\citeproctext]{ref-lawson_speedy_2019}
\textbf{\textbf{Lawson T}, \textbf{Vialet-Chabrand S}}. \textbf{2019}.
\href{https://doi.org/10.1111/nph.15330}{Speedy stomata, photosynthesis
and plant water use efficiency}. \emph{New Phytologist} \textbf{221}:
93--98.

\bibitem[\citeproctext]{ref-lawson_photosynthesis_2010}
\textbf{\textbf{Lawson T}, \textbf{Von Caemmerer S}, \textbf{Baroli I}}.
\textbf{2010}.
\href{https://doi.org/10.1007/978-3-642-13145-5_11}{Photosynthesis and
{Stomatal Behaviour}}. In: Lüttge UE, Beyschlag W, Büdel B, Francis D,
eds. Progress in {Botany} 72. Berlin, Heidelberg: Springer Berlin
Heidelberg, 265--304.

\bibitem[\citeproctext]{ref-lehmann_effects_2015}
\textbf{\textbf{Lehmann P}, \textbf{Or D}}. \textbf{2015}.
\href{https://doi.org/10.1111/nph.13442}{Effects of stomata clustering
on leaf gas exchange}. \emph{New Phytologist} \textbf{207}: 1015--1025.

\bibitem[\citeproctext]{ref-lynch_methods_1991}
\textbf{\textbf{Lynch M}}. \textbf{1991}.
\href{https://doi.org/10.1111/j.1558-5646.1991.tb04375.x}{Methods for
the analysis of comparative data in evolutionary biology}.
\emph{Evolution} \textbf{45}: 1065--1080.

\bibitem[\citeproctext]{ref-matthews_role_2020}
\textbf{\textbf{Matthews JSA}, \textbf{Vialet-Chabrand S},
\textbf{Lawson T}}. \textbf{2020}.
\href{https://doi.org/10.1093/jxb/erz563}{Role of blue and red light in
stomatal dynamic behaviour}. \emph{Journal of Experimental Botany}
\textbf{71}: 2253--2269.

\bibitem[\citeproctext]{ref-mcadam_stomatal_2012}
\textbf{\textbf{McAdam SAM}, \textbf{Brodribb TJ}}. \textbf{2012}.
\href{https://doi.org/10.1111/j.1461-0248.2011.01700.x}{Stomatal
innovation and the rise of seed plants}. \emph{Ecology Letters}
\textbf{15}: 1--8.

\bibitem[\citeproctext]{ref-mcadam_evolution_2015}
\textbf{\textbf{McAdam SAM}, \textbf{Brodribb TJ}}. \textbf{2015}.
\href{https://doi.org/10.1104/pp.114.252940}{The {Evolution} of
{Mechanisms Driving} the {Stomatal Response} to {Vapor Pressure
Deficit}}. \emph{Plant Physiology} \textbf{167}: 833--843.

\bibitem[\citeproctext]{ref-mcadam_stomatal_2016}
\textbf{\textbf{McAdam SAM}, \textbf{Sussmilch FC}, \textbf{Brodribb
TJ}}. \textbf{2016}. \href{https://doi.org/10.1111/pce.12633}{Stomatal
responses to vapour pressure deficit are regulated by high speed gene
expression in angiosperms}. \emph{Plant, Cell \& Environment}
\textbf{39}: 485--491.

\bibitem[\citeproctext]{ref-mcausland_effects_2016}
\textbf{\textbf{McAusland L}, \textbf{Vialet-Chabrand S}, \textbf{Davey
P}, \textbf{Baker NR}, \textbf{Brendel O}, \textbf{Lawson T}}.
\textbf{2016}. \href{https://doi.org/10.1111/nph.14000}{Effects of
kinetics of light-induced stomatal responses on photosynthesis and
water-use efficiency}. \emph{New Phytologist} \textbf{211}: 1209--1220.

\bibitem[\citeproctext]{ref-mcelreath_statistical_2020}
\textbf{\textbf{McElreath R}}. \textbf{2020}. \emph{Statistical
{Rethinking}: A {Bayesian} course with examples in {R} and {Stan}}. Boca
Raton: {Taylor and Francis, CRC Press}.

\bibitem[\citeproctext]{ref-mcelwain_using_2016}
\textbf{\textbf{McElwain JC}, \textbf{Yiotis C}, \textbf{Lawson T}}.
\textbf{2016}. \href{https://doi.org/10.1111/nph.13579}{Using modern
plant trait relationships between observed and theoretical maximum
stomatal conductance and vein density to examine patterns of plant
macroevolution}. \emph{New Phytologist} \textbf{209}: 94--103.

\bibitem[\citeproctext]{ref-mclatchie_efficient_2024}
\textbf{\textbf{McLatchie Y}, \textbf{Vehtari A}}. \textbf{2024}.
\href{https://doi.org/10.1007/s11222-024-10442-4}{Efficient estimation
and correction of selection-induced bias with order statistics}.
\emph{Statistics and Computing} \textbf{34}: 132.

\bibitem[\citeproctext]{ref-medlyn_reconciling_2011}
\textbf{\textbf{Medlyn BE}, \textbf{Duursma RA}, \textbf{Eamus D},
\textbf{Ellsworth DS}, \textbf{Prentice IC}, \textbf{Barton CVM},
\textbf{Crous KY}, \textbf{De Angelis P}, \textbf{Freeman M},
\textbf{Wingate L}}. \textbf{2011}.
\href{https://doi.org/10.1111/j.1365-2486.2010.02375.x}{Reconciling the
optimal and empirical approaches to modelling stomatal conductance}.
\emph{Global Change Biology} \textbf{17}: 2134--2144.

\bibitem[\citeproctext]{ref-merilo_open_2014}
\textbf{\textbf{Merilo E}, \textbf{Jõesaar I}, \textbf{Brosché M},
\textbf{Kollist H}}. \textbf{2014}.
\href{https://doi.org/10.1111/nph.12667}{To open or to close:
Species-specific stomatal responses to simultaneously applied opposing
environmental factors}. \emph{New Phytologist} \textbf{202}: 499--508.

\bibitem[\citeproctext]{ref-monteith_reinterpretation_1995}
\textbf{\textbf{Monteith JL}}. \textbf{1995}.
\href{https://doi.org/10.1111/j.1365-3040.1995.tb00371.x}{A
reinterpretation of stomatal responses to humidity}. \emph{Plant, Cell
\& Environment} \textbf{18}: 357--364.

\bibitem[\citeproctext]{ref-monteith_principles_2013}
\textbf{\textbf{Monteith JL}, \textbf{Unsworth MH}}. \textbf{2013}.
\emph{Principles of environmental physics: Plants, animals, and the
atmosphere}. Amsterdam ; Boston: Elsevier/Academic Press.

\bibitem[\citeproctext]{ref-mott_stomatal_1991}
\textbf{\textbf{Mott KA}, \textbf{Parkhurst DF}}. \textbf{1991}.
\href{https://doi.org/10.1111/j.1365-3040.1991.tb01521.x}{Stomatal
responses to humidity in air and helox}. \emph{Plant, Cell \&
Environment} \textbf{14}: 509--515.

\bibitem[\citeproctext]{ref-muir_how_2023}
\textbf{\textbf{Muir CD}, \textbf{Conesa MÀ}, \textbf{Galmés J},
\textbf{Pathare VS}, \textbf{Rivera P}, \textbf{López Rodríguez R},
\textbf{Terrazas T}, \textbf{Xiong D}}. \textbf{2023}.
\href{https://doi.org/10.1086/723780}{How important are functional and
developmental constraints on phenotypic evolution? {An} empirical test
with the stomatal anatomy of flowering plants}. \emph{The American
Naturalist} \textbf{201}: 794--812.

\bibitem[\citeproctext]{ref-muir_plasticity_2025}
\textbf{\textbf{Muir CD}, \textbf{Lim WS}, \textbf{Wang D}}.
\textbf{2025}.
\href{https://doi.org/10.1101/2025.08.19.671015}{Plasticity and
adaptation to high light intensity amplify the advantage of
amphistomatous leaves}.

\bibitem[\citeproctext]{ref-murray_consistent_2020}
\textbf{\textbf{Murray M}, \textbf{Soh WK}, \textbf{Yiotis C},
\textbf{Spicer RA}, \textbf{Lawson T}, \textbf{McElwain JC}}.
\textbf{2020}. \href{https://doi.org/10.1086/706260}{Consistent
relationship between field-measured stomatal conductance and theoretical
maximum stomatal conductance in {C}{\textsubscript{3}} woody angiosperms
in four major biomes}. \emph{International Journal of Plant Sciences}
\textbf{181}: 142--154.

\bibitem[\citeproctext]{ref-ochoa_pinpointing_2024}
\textbf{\textbf{Ochoa ME}, \textbf{Henry C}, \textbf{John GP},
\textbf{Medeiros CD}, \textbf{Pan R}, \textbf{Scoffoni C},
\textbf{Buckley TN}, \textbf{Sack L}}. \textbf{2024}.
\href{https://doi.org/10.1093/plphys/kiae292}{Pinpointing the causal
influences of stomatal anatomy and behavior on minimum, operational, and
maximum leaf surface conductance}. \emph{Plant Physiology} \textbf{196}:
51--66.

\bibitem[\citeproctext]{ref-papanatsiou_optogenetic_2019}
\textbf{\textbf{Papanatsiou M}, \textbf{Petersen J}, \textbf{Henderson
L}, \textbf{Wang Y}, \textbf{Christie JM}, \textbf{Blatt MR}}.
\textbf{2019}.
\href{https://doi.org/10.1126/science.aaw0046}{Optogenetic manipulation
of stomatal kinetics improves carbon assimilation, water use, and
growth}. \emph{Science} \textbf{363}: 1456--1459.

\bibitem[\citeproctext]{ref-pearl_causality_2022}
\textbf{\textbf{Pearl J}}. \textbf{2022}. \emph{Causality: Models,
reasoning, and inference}. Cambridge New York, NY Port Melbourne New
Delhi Singapore: Cambridge University Press.

\bibitem[\citeproctext]{ref-pease_phylogenomics_2016}
\textbf{\textbf{Pease JB}, \textbf{Haak DC}, \textbf{Hahn MW},
\textbf{Moyle LC}}. \textbf{2016}.
\href{https://doi.org/10.1371/journal.pbio.1002379}{Phylogenomics
reveals three sources of adaptive variation during a rapid radiation}.
\emph{PLOS Biology} \textbf{14}: e1002379.

\bibitem[\citeproctext]{ref-peralta_taxonomy_2008}
\textbf{\textbf{Peralta IE}, \textbf{Spooner DM}, \textbf{Knapp S}}.
\textbf{2008}. Taxonomy of wild tomatoes and their relatives
({\emph{Solanum}} sect. {\emph{Lycopersicoides}}, sect.
{\emph{Juglandifolia}}, sect. {\emph{Lycopersicon}}; {Solanaceae}).
\textbf{84}.

\bibitem[\citeproctext]{ref-pichaco_stomataltoepidermal_2025}
\textbf{\textbf{Pichaco J}, \textbf{Diaz-Espejo A},
\textbf{Rodriguez-Dominguez CM}}. \textbf{2025}.
\href{https://doi.org/10.1111/ppl.70543}{Stomatal-{To-Epidermal Cell
Size Ratio} and {Physiological Coordination Influence Stomatal Response
Speed} to {Leaf Water Disruption} in {Tomato Plants}}. \emph{Physiologia
Plantarum} \textbf{177}: e70543.

\bibitem[\citeproctext]{ref-pichaco_mechanical_2024}
\textbf{\textbf{Pichaco J}, \textbf{Manandhar A}, \textbf{McAdam SAM}}.
\textbf{2024}. \href{https://doi.org/10.1093/plphys/kiae023}{Mechanical
advantage makes stomatal opening speed a function of evaporative
demand}. \emph{Plant Physiology} \textbf{195}: 370--377.

\bibitem[\citeproctext]{ref-pichaco_hydroactive_2026}
\textbf{\textbf{Pichaco J}, \textbf{Rodriguez-Dominguez CM},
\textbf{Diaz-Espejo A}}. \textbf{2026}.
\href{https://doi.org/10.1016/j.tplants.2026.07.013}{Hydroactive
stomata: A key innovation underpinning angiosperm ecological dominance}.
\emph{Trends in Plant Science}: S1360138526002438.

\bibitem[\citeproctext]{ref-potkay_generalized_2025}
\textbf{\textbf{Potkay A}, \textbf{Cabon A}, \textbf{Peters RL},
\textbf{Fonti P}, \textbf{Sapes G}, \textbf{Sala A}, \textbf{Stefanski
A}, \textbf{Butler E}, \textbf{Bermudez R}, \textbf{Montgomery R},
\emph{et al.}} \textbf{2025}.
\href{https://doi.org/10.1111/gcb.70049}{Generalized {Stomatal
Optimization} of {Evolutionary Fitness Proxies} for {Predicting Plant
Gas Exchange Under Drought}, {Heatwaves}, and {Elevated
CO}{\textsubscript{2}}}. \emph{Global Change Biology} \textbf{31}:
e70049.

\bibitem[\citeproctext]{ref-prentice_balancing_2014}
\textbf{\textbf{Prentice IC}, \textbf{Dong N}, \textbf{Gleason SM},
\textbf{Maire V}, \textbf{Wright IJ}}. \textbf{2014}.
\href{https://doi.org/10.1111/ele.12211}{Balancing the costs of carbon
gain and water transport: Testing a new theoretical framework for plant
functional ecology}. \emph{Ecology Letters} \textbf{17}: 82--91.

\bibitem[\citeproctext]{ref-qie_stomatal_2024}
\textbf{\textbf{Qie Y-D}, \textbf{Zhang Q-W}, \textbf{McAdam SAM},
\textbf{Cao K-F}}. \textbf{2024}.
\href{https://doi.org/10.1016/j.pld.2024.02.003}{Stomatal dynamics are
regulated by leaf hydraulic traits and guard cell anatomy in nine true
mangrove species}. \emph{Plant Diversity} \textbf{46}: 395--405.

\bibitem[\citeproctext]{ref-rcoreteam_language_2026}
\textbf{\textbf{R Core Team}}. \textbf{2026}. \emph{R: {A Language} and
{Environment} for {Statistical Computing}}. Vienna, Austria: R
Foundation for Statistical Computing.

\bibitem[\citeproctext]{ref-raven_selection_2002}
\textbf{\textbf{Raven JA}}. \textbf{2002}.
\href{https://doi.org/10.1046/j.0028-646X.2001.00334.x}{Selection
pressures on stomatal evolution}. \emph{New Phytologist} \textbf{153}:
371--386.

\bibitem[\citeproctext]{ref-raven_speedy_2014}
\textbf{\textbf{Raven JA}}. \textbf{2014}.
\href{https://doi.org/10.1093/jxb/eru032}{Speedy small stomata?}
\emph{Journal of Experimental Botany} \textbf{65}: 1415--1424.

\bibitem[\citeproctext]{ref-rockwell_extreme_2022}
\textbf{\textbf{Rockwell FE}, \textbf{Holbrook NM}, \textbf{Jain P},
\textbf{Huber AE}, \textbf{Sen S}, \textbf{Stroock AD}}. \textbf{2022}.
\href{https://doi.org/10.1093/aob/mcac094}{Extreme undersaturation in
the intercellular airspace of leaves: A failure of {Gaastra} or {Ohm}?}
\emph{Annals of Botany} \textbf{130}: 301--316.

\bibitem[\citeproctext]{ref-sack_developmental_2016}
\textbf{\textbf{Sack L}, \textbf{Buckley TN}}. \textbf{2016}.
\href{https://doi.org/10.1104/pp.16.00476}{The developmental basis of
stomatal density and flux}. \emph{Plant Physiology} \textbf{171}:
2358--2363.

\bibitem[\citeproctext]{ref-schoch_dependence_1980}
\textbf{\textbf{Schoch P-G}, \textbf{Zinsou C}, \textbf{Sibi M}}.
\textbf{1980}. \href{https://doi.org/10.1093/jxb/31.5.1211}{Dependence
of the stomatal index on environmental factors during stomatal
differentiation in leaves of {\emph{Vigna}}{ \emph{Sinensis}} {L}.: 1.
{Effect} of light intensity}. \emph{Journal of Experimental Botany}
\textbf{31}: 1211--1216.

\bibitem[\citeproctext]{ref-schymanski_stomatal_2013}
\textbf{\textbf{Schymanski SJ}, \textbf{Or D}, \textbf{Zwieniecki M}}.
\textbf{2013}.
\href{https://doi.org/10.1371/journal.pone.0054231}{Stomatal control and
leaf thermal and hydraulic capacitances under rapid environmental
fluctuations}. \emph{PLoS ONE} \textbf{8}: e54231.

\bibitem[\citeproctext]{ref-smaldino_natural_2016}
\textbf{\textbf{Smaldino PE}, \textbf{McElreath R}}. \textbf{2016}.
\href{https://doi.org/10.1098/rsos.160384}{The natural selection of bad
science}. \emph{Royal Society Open Science} \textbf{3}: 160384.

\bibitem[\citeproctext]{ref-sperry_predicting_2017}
\textbf{\textbf{Sperry JS}, \textbf{Venturas MD}, \textbf{Anderegg WRL},
\textbf{Mencuccini M}, \textbf{Mackay DS}, \textbf{Wang Y}, \textbf{Love
DM}}. \textbf{2017}. \href{https://doi.org/10.1111/pce.12852}{Predicting
stomatal responses to the environment from the optimization of
photosynthetic gain and hydraulic cost: {A} stomatal optimization
model}. \emph{Plant, Cell \& Environment} \textbf{40}: 816--830.

\bibitem[\citeproctext]{ref-standevelopmentteam_stan_2026}
\textbf{\textbf{Stan Development Team}}. \textbf{2026}. \emph{Stan
{Modeling Language Users Guide} and {Reference Manual}}.

\bibitem[\citeproctext]{ref-tang_variation_2026}
\textbf{\textbf{Tang X}, \textbf{Huang W}, \textbf{Yuan X}, \textbf{Liu
X}, \textbf{Wan J}}. \textbf{2026}.
\href{https://doi.org/10.1111/ppl.70824}{Variation in {Stomatal Traits}
and {Photosynthetic Induction Between Subtropical Evergreen} and
{Deciduous Broadleaved Tree Species}}. \emph{Physiologia Plantarum}
\textbf{178}: e70824.

\bibitem[\citeproctext]{ref-tezara_water_1999}
\textbf{\textbf{Tezara W}, \textbf{Mitchell VJ}, \textbf{Driscoll SD},
\textbf{Lawlor DW}}. \textbf{1999}.
\href{https://doi.org/10.1038/44842}{Water stress inhibits plant
photosynthesis by decreasing coupling factor and {ATP}}. \emph{Nature}
\textbf{401}: 914--917.

\bibitem[\citeproctext]{ref-ting_further_1965}
\textbf{\textbf{Ting IP}, \textbf{Loomis WE}}. \textbf{1965}.
\href{https://doi.org/10.1104/pp.40.2.220}{Further {Studies Concerning
Stomatal Diffusion}}. \emph{Plant Physiology} \textbf{40}: 220--228.

\bibitem[\citeproctext]{ref-uyeda_rethinking_2018}
\textbf{\textbf{Uyeda JC}, \textbf{Zenil-Ferguson R}, \textbf{Pennell
MW}}. \textbf{2018}.
\href{https://doi.org/10.1093/sysbio/syy031}{Rethinking phylogenetic
comparative methods} (N Matzke, Ed.). \emph{Systematic Biology}
\textbf{67}: 1091--1109.

\bibitem[\citeproctext]{ref-vehtari_practical_2017}
\textbf{\textbf{Vehtari A}, \textbf{Gelman A}, \textbf{Gabry J}}.
\textbf{2017}.
\href{https://doi.org/10.1007/s11222-016-9696-4}{Practical {Bayesian}
model evaluation using leave-one-out cross-validation and {WAIC}}.
\emph{Statistics and Computing} \textbf{27}: 1413--1432.

\bibitem[\citeproctext]{ref-vico_effects_2011}
\textbf{\textbf{Vico G}, \textbf{Manzoni S}, \textbf{Palmroth S},
\textbf{Katul G}}. \textbf{2011}.
\href{https://doi.org/10.1111/j.1469-8137.2011.03847.x}{Effects of
stomatal delays on the economics of leaf gas exchange under intermittent
light regimes}. \emph{New Phytologist} \textbf{192}: 640--652.

\bibitem[\citeproctext]{ref-wall_stomata_2022}
\textbf{\textbf{Wall S}, \textbf{Vialet-Chabrand S}, \textbf{Davey P},
\textbf{Van Rie J}, \textbf{Galle A}, \textbf{Cockram J}, \textbf{Lawson
T}}. \textbf{2022}. \href{https://doi.org/10.1111/nph.18257}{Stomata on
the abaxial and adaxial leaf surfaces contribute differently to leaf gas
exchange and photosynthesis in wheat}. \emph{New Phytologist}
\textbf{235}: 1743--1756.

\bibitem[\citeproctext]{ref-westbrook_stomatal_2021}
\textbf{\textbf{Westbrook AS}, \textbf{McAdam SAM}}. \textbf{2021}.
\href{https://doi.org/10.1111/nph.16850}{Stomatal density and mechanics
are critical for high productivity: Insights from amphibious ferns}.
\emph{New Phytologist} \textbf{229}: 877--889.

\bibitem[\citeproctext]{ref-westoby_misinterpreting_1995}
\textbf{\textbf{Westoby M}, \textbf{Leishman MR}, \textbf{Lord JM}}.
\textbf{1995}. \href{https://doi.org/10.2307/2261605}{On
{Misinterpreting} the {`{Phylogenetic Correction}'}}. \emph{The Journal
of Ecology} \textbf{83}: 531.

\bibitem[\citeproctext]{ref-westoby_phylogenetically_2023}
\textbf{\textbf{Westoby M}, \textbf{Yates L}, \textbf{Holland B},
\textbf{Halliwell B}}. \textbf{2023}.
\href{https://doi.org/10.1111/1365-2745.14150}{Phylogenetically
conservative trait correlation: {Quantification} and interpretation}.
\emph{Journal of Ecology} \textbf{111}: 2105--2117.

\bibitem[\citeproctext]{ref-willmer_stomata_1996}
\textbf{\textbf{Willmer C}, \textbf{Fricker M}}. \textbf{1996}.
\emph{\href{https://doi.org/10.1007/978-94-011-0579-8}{Stomata}}.
Dordrecht: Springer Netherlands.

\bibitem[\citeproctext]{ref-wong_humidity_2022}
\textbf{\textbf{Wong SC}, \textbf{Canny MJ}, \textbf{Holloway-Phillips
M}, \textbf{Stuart-Williams H}, \textbf{Cernusak LA}, \textbf{Márquez
DA}, \textbf{Farquhar GD}}. \textbf{2022}.
\href{https://doi.org/10.1038/s41477-022-01202-1}{Humidity gradients in
the air spaces of leaves}. \emph{Nature Plants} \textbf{8}: 971--978.

\bibitem[\citeproctext]{ref-woning_revisiting_2026}
\textbf{\textbf{Woning N}, \textbf{Al-Salman Y}, \textbf{Kaiser E},
\textbf{Berman SR}, \textbf{Brendel O}, \textbf{Cano FJ},
\textbf{Carpentier S}, \textbf{Centritto M}, \textbf{Drake PL},
\textbf{Durand M}, \emph{et al.}} \textbf{2026}.
\href{https://doi.org/10.1111/nph.70842}{Revisiting the relationship
between stomatal size and speed across species -- a meta-analysis}.
\emph{New Phytologist} \textbf{249}: 2338--2354.

\bibitem[\citeproctext]{ref-xie_stomatal_2022}
\textbf{\textbf{Xie J}, \textbf{Wang Z}, \textbf{Li Y}}. \textbf{2022}.
\href{https://doi.org/10.1111/nph.18189}{Stomatal opening ratio mediates
trait coordinating network adaptation to environmental gradients}.
\emph{New Phytologist} \textbf{235}: 907--922.

\bibitem[\citeproctext]{ref-xie_identification_2006}
\textbf{\textbf{Xie X}, \textbf{Wang Y}, \textbf{Williamson L},
\textbf{Holroyd GH}, \textbf{Tagliavia C}, \textbf{Murchie E},
\textbf{Theobald J}, \textbf{Knight MR}, \textbf{Davies WJ},
\textbf{Leyser HMO}, \emph{et al.}} \textbf{2006}.
\href{https://doi.org/10.1016/j.cub.2006.03.028}{The {Identification} of
{Genes Involved} in the {Stomatal Response} to {Reduced Atmospheric
Relative Humidity}}. \emph{Current Biology} \textbf{16}: 882--887.

\bibitem[\citeproctext]{ref-yoshiyama_natural_2024}
\textbf{\textbf{Yoshiyama Y}, \textbf{Wakabayashi Y}, \textbf{Mercer
KL}, \textbf{Kawabata S}, \textbf{Kobayashi T}, \textbf{Tabuchi T},
\textbf{Yamori W}}. \textbf{2024}.
\href{https://doi.org/10.1093/jxb/erae082}{Natural genetic variation in
dynamic photosynthesis is correlated with stomatal anatomical traits in
diverse tomato species across geographical habitats} (T Lawson, Ed.).
\emph{Journal of Experimental Botany} \textbf{75}: 6762--6777.

\end{CSLReferences}

\newpage{}

\subsection{ALT TEXT figures}\label{alt-text-figures}

\textbf{Figure 1.} Four-panel conceptual diagram. Panel a: stomatal
conductance over time after a step increase in vapor pressure deficit,
showing a brief rise (wrong-way response) followed by a decline to a
lower steady state (right-way response); an orange curve shows a Weibull
function fit to the decline, with initial and final conductance values
marked. Panel b: four decline curves showing how the time-constant and
lag-time parameters change the response shape. The smooth, roughly
exponential declines, have low lag time at different rates; the two
S-shaped curves show accelerating decline at different rates after a
lag. Panel c: Two curves showing faster closure kinetics in smaller
guard cells with greater surface area surface-area-to-volume ratio.
Panel d: a saturating curve of stomatal conductance versus guard cell
turgor pressure, with tangent lines at two points showing the slope is
steeper at low conductance (low fraction of maximum conductance) and
shallower at high conductance (high fraction of maximum conductance).

\textbf{Figure 2.} Horizontal stacked bar chart showing the percentage
of variance apportioned among three hierarchical levels ---
phylogenetic, non-phylogenetic population, and among-individual --- for
five traits: guard cell length, anatomical maximum stomatal conductance,
the stomatal closure time constant, initial stomatal conductance, and
the lag time. Guard cell length, maximum conductance, and the time
constant have long phylogenetic segments, whereas initial conductance
and lag time have long among-individual segments; the non-phylogenetic
population segment is comparatively short for every trait.

\textbf{Figure 3.} Two-panel figure of population-level trait means
across experimental treatments, with points for each wild tomato
population connected by thin gray lines across treatment levels. Panel
a: guard cell length values arranged by growth light intensity (shade,
sun) and leaf type. Panel b: the fraction of anatomical maximum stomatal
conductance arranged by measurement light intensity and leaf type.

\textbf{Figure 4.} Faceted dot plot arranged in a two-by-two grid. Rows
show the two kinetic parameters, lag time (top) and time constant
(bottom); columns show the growth light treatment (shade, left; sun,
right). Within each panel, points represent the average value for one
wild tomato population under low or high measurement light intensity,
allowing comparison of parameter values across light intensities and
growth treatments.

\textbf{Figure 5.} Scatter plot of guard cell length (horizontal axis,
log scale) against the stomatal closure time constant (vertical axis,
log scale) for wild tomato populations, arranged in a two-by-two grid of
facets by growth light intensity (columns) and measurement light
intensity (rows). Points are colored by leaf type (amphistomatous
vs.~pseudohypostomatous) and overlaid with colored ellipses representing
the phylogenetic covariance between the two traits in each treatment
combination; the ellipses are tilted with a positive slope in every
facet.

\textbf{Figure 6.} Scatter plot of initial stomatal conductance
(horizontal axis) against the stomatal closure time constant (vertical
axis) for individual leaves, arranged in a two-by-two grid of facets by
growth light intensity (columns) and measurement light intensity (rows).
Points are colored by leaf type (amphistomatous
vs.~pseudohypostomatous), each with a fitted straight regression line
and shaded confidence ribbon showing a positive relationship between
initial conductance and the time constant in every facet.

\textbf{Figure 7.} Path diagram with three stacked panels, one for each
experimental treatment: growth light intensity (top), leaf type
(middle), and measurement light intensity (bottom). In each panel, a
horizontal arrow represents the direct effect of the treatment on the
stomatal closure time constant, and a pair of slanted arrows passing
through a node for initial stomatal conductance represents the indirect
effect. Arrow thickness is scaled to effect size, and arrows are colored
red for negative effects or blue for positive effects, with numeric
effect estimates and confidence intervals labeled above each arrow.

\newpage{}

\subsection{Supporting Information
legends}\label{supporting-information-legends}

\begin{flushleft}

\textbf{Table S1.} Accession information of \emph{Solanum} populations
used in this study.

\textbf{Table S2.} Relative humidity and vapor pressure deficit during
stomatal kinetics measurements.

\textbf{Table S3.} Total number of humidity response curves analyzed by
treatment combination.

\textbf{Table S4.} Model comparison results for fixed effect of guard
cell length and \fgmax{} components on kinetic parameters.

\textbf{Table S5.} Leave-one-out cross-validation information criteria
(LOOIC) for plausible models.

\textbf{Table S6.} Relative contributions of phylogenetic, among
population (nonphylogenetic), and among individual varation among
stomatal traits.

\textbf{Table S7.} Fixed effect parameter estimates from the selected
Bayesian model.

\textbf{Table S8.} Estimates of stomatal anatomy and kinetic variables
associated with each curve.

\textbf{Table S9.} Estimates of stomatal anatomy and kinetic variables
associated with each population and treatment combination.

\textbf{Table S10.} Random effect parameter estimates from the selected
Bayesian model.

\textbf{Figure S1.} Pairwise relationships among vapor pressure deficit
(VPD) covariates.

\textbf{Figure S2.} Correlation between stomatal closure time constant
\(\tau\) and the inverse VPD saturation rate.

\textbf{Figure S3.} Vapor pressure deficit saturation rate in sample and
references airstreams.

\textbf{Figure S4.} Estimated effects of \(g_\mathrm{i}\) on \(\tau\)
and phylogenetic correlation between guard cell length and \(\tau\) in
selected model and VPD-adjusted model.

\textbf{Figure S5.} Estimated and observed maximum stomatal
conductances.

\textbf{Figure S6.} Humidity-response curves and fitted lines.

\textbf{Figure S7.} Comparison of adaxial and abaxial guard cell lengths
among populations.

\textbf{Figure S8.} Individual-level effects across plausible models.

\textbf{Figure S9.} Phylogenetic correlations across plausible models.

\textbf{Figure S10.} Posterior estimates of phylogenetic correlations
and fixed effects of guard cell length on \(\tau\).

\textbf{Figure S11.} Individual-level variation in maximum conductance
and the time constant \(\tau\).

\textbf{Figure S12.} Individual-level variation in initial conductance
influences lag time \(\lambda\).

\textbf{Figure S13.} Individual-level variation in the anatomical
maximum conductance does not influence lag time \(\lambda\).

\textbf{Notes S1.} Effects of \(g_\mathrm{i}\) and guard cell length are
robust to variation in VPD stimulus.

\textbf{Notes S2.} The \(g_\mathrm{i}\)--\(\tau\) association is not an
artifact of the curve-fitting procedure.

\textbf{Notes S3.} Sensitivity analysis of the
\(g_\mathrm{i}\)--\(\tau\) path analysis to unmeasured confounding.

\textbf{Figure A.} Comparison of the observed \(g_\mathrm{i}\)--\(\tau\)
correlation against a null distribution.

\textbf{Figure B.} Causal structure of the \(g_\mathrm{i}\)--\(\tau\)
path analysis and unmeasured confounding.

\textbf{Figure C.} Sensitivity of indirect treatment effects on \(\tau\)
through \(g_\mathrm{i}\) to residual correlation assumptions.

\end{flushleft}

\end{document}
